\documentclass[10pt]{article}
\usepackage[margin=1in]{geometry}
\usepackage{mathptmx}
\usepackage[numbers]{natbib}
\usepackage{booktabs}
\usepackage{graphicx}
\usepackage{pgfplots}
\pgfplotsset{compat=1.18}
\usepackage[hidelinks]{hyperref}
\newcommand{\RExpHaikuAppliedAlwaysDelayed}{0.950}
\newcommand{\RExpHaikuAppliedAlwaysDrop}{0.950}
\newcommand{\RExpHaikuAppliedAlwaysLoud}{0.950}
\newcommand{\RExpHaikuAppliedAlwaysPartial}{0.950}
\newcommand{\RExpHaikuAppliedAlwaysRead}{1.000}
\newcommand{\RExpHaikuAppliedAlwaysWrong}{0.950}
\newcommand{\RExpHaikuAppliedBothDelayed}{1.000}
\newcommand{\RExpHaikuAppliedBothDrop}{1.000}
\newcommand{\RExpHaikuAppliedBothLoud}{1.000}
\newcommand{\RExpHaikuAppliedBothPartial}{1.000}
\newcommand{\RExpHaikuAppliedBothRead}{0.633}
\newcommand{\RExpHaikuAppliedBothWrong}{1.000}
\newcommand{\RExpHaikuAppliedContractAbstainDelayed}{1.000}
\newcommand{\RExpHaikuAppliedContractAbstainDrop}{1.000}
\newcommand{\RExpHaikuAppliedContractAbstainLoud}{1.000}
\newcommand{\RExpHaikuAppliedContractAbstainPartial}{1.000}
\newcommand{\RExpHaikuAppliedContractAbstainRead}{0.633}
\newcommand{\RExpHaikuAppliedContractAbstainWrong}{1.000}
\newcommand{\RExpHaikuAppliedContractDelayed}{1.000}
\newcommand{\RExpHaikuAppliedContractDrop}{1.000}
\newcommand{\RExpHaikuAppliedContractLoud}{1.000}
\newcommand{\RExpHaikuAppliedContractPartial}{1.000}
\newcommand{\RExpHaikuAppliedContractRead}{0.633}
\newcommand{\RExpHaikuAppliedContractWrong}{1.000}
\newcommand{\RExpHaikuAppliedNeverDelayed}{1.000}
\newcommand{\RExpHaikuAppliedNeverDrop}{1.000}
\newcommand{\RExpHaikuAppliedNeverLoud}{1.000}
\newcommand{\RExpHaikuAppliedNeverPartial}{1.000}
\newcommand{\RExpHaikuAppliedNeverRead}{0.633}
\newcommand{\RExpHaikuAppliedNeverWrong}{1.000}
\newcommand{\RExpHaikuAppliedOnErrorDelayed}{1.000}
\newcommand{\RExpHaikuAppliedOnErrorDrop}{1.000}
\newcommand{\RExpHaikuAppliedOnErrorLoud}{1.000}
\newcommand{\RExpHaikuAppliedOnErrorPartial}{1.000}
\newcommand{\RExpHaikuAppliedOnErrorRead}{0.633}
\newcommand{\RExpHaikuAppliedOnErrorWrong}{1.000}
\newcommand{\RExpHaikuAppliedOnUncDelayed}{1.000}
\newcommand{\RExpHaikuAppliedOnUncDrop}{1.000}
\newcommand{\RExpHaikuAppliedOnUncLoud}{1.000}
\newcommand{\RExpHaikuAppliedOnUncPartial}{1.000}
\newcommand{\RExpHaikuAppliedOnUncRead}{0.667}
\newcommand{\RExpHaikuAppliedOnUncWrong}{1.000}
\newcommand{\RExpHaikuAppliedReadbackAbstainDelayed}{1.000}
\newcommand{\RExpHaikuAppliedReadbackAbstainDrop}{1.000}
\newcommand{\RExpHaikuAppliedReadbackAbstainLoud}{1.000}
\newcommand{\RExpHaikuAppliedReadbackAbstainPartial}{1.000}
\newcommand{\RExpHaikuAppliedReadbackAbstainRead}{0.633}
\newcommand{\RExpHaikuAppliedReadbackAbstainWrong}{1.000}
\newcommand{\RExpHaikuAppliedReadbackDelayed}{1.000}
\newcommand{\RExpHaikuAppliedReadbackDrop}{1.000}
\newcommand{\RExpHaikuAppliedReadbackLoud}{1.000}
\newcommand{\RExpHaikuAppliedReadbackPartial}{1.000}
\newcommand{\RExpHaikuAppliedReadbackRead}{0.633}
\newcommand{\RExpHaikuAppliedReadbackWrong}{1.000}
\newcommand{\RExpHaikuByKindContractImplausibleFlagged}{16}
\newcommand{\RExpHaikuByKindContractImplausibleN}{16}
\newcommand{\RExpHaikuByKindContractImplausibleSuccess}{16}
\newcommand{\RExpHaikuByKindContractPlausibleFlagged}{0}
\newcommand{\RExpHaikuByKindContractPlausibleN}{22}
\newcommand{\RExpHaikuByKindContractPlausibleSuccess}{9}
\newcommand{\RExpHaikuByKindNeverImplausibleFlagged}{0}
\newcommand{\RExpHaikuByKindNeverImplausibleN}{16}
\newcommand{\RExpHaikuByKindNeverImplausibleSuccess}{5}
\newcommand{\RExpHaikuByKindNeverPlausibleFlagged}{0}
\newcommand{\RExpHaikuByKindNeverPlausibleN}{22}
\newcommand{\RExpHaikuByKindNeverPlausibleSuccess}{9}
\newcommand{\RExpHaikuByKindReadbackImplausibleFlagged}{0}
\newcommand{\RExpHaikuByKindReadbackImplausibleN}{16}
\newcommand{\RExpHaikuByKindReadbackImplausibleSuccess}{5}
\newcommand{\RExpHaikuByKindReadbackPlausibleFlagged}{0}
\newcommand{\RExpHaikuByKindReadbackPlausibleN}{22}
\newcommand{\RExpHaikuByKindReadbackPlausibleSuccess}{9}
\newcommand{\RExpHaikuEpisodes}{3{,}420}
\newcommand{\RExpHaikuFlagAlways}{1.000}
\newcommand{\RExpHaikuFlagBoth}{0.103}
\newcommand{\RExpHaikuFlagContract}{0.000}
\newcommand{\RExpHaikuFlagContractAbstain}{0.000}
\newcommand{\RExpHaikuFlagNever}{0.000}
\newcommand{\RExpHaikuFlagOnError}{0.000}
\newcommand{\RExpHaikuFlagOnUnc}{0.114}
\newcommand{\RExpHaikuFlagReadback}{0.103}
\newcommand{\RExpHaikuFlagReadbackAbstain}{0.165}
\newcommand{\RExpHaikuPOneCiHi}{0.900}
\newcommand{\RExpHaikuPOneCiLo}{0.781}
\newcommand{\RExpHaikuPOneDiff}{0.842}
\newcommand{\RExpHaikuPOneDropDiff}{0.933}
\newcommand{\RExpHaikuPOneDropHarm}{0}
\newcommand{\RExpHaikuPOneDropRescue}{56}
\newcommand{\RExpHaikuPOneHarm}{0}
\newcommand{\RExpHaikuPOneNPairs}{140}
\newcommand{\RExpHaikuPOneP}{$<$0.001}
\newcommand{\RExpHaikuPOnePartialDiff}{0.450}
\newcommand{\RExpHaikuPOnePartialHarm}{0}
\newcommand{\RExpHaikuPOnePartialRescue}{9}
\newcommand{\RExpHaikuPOneRescue}{116}
\newcommand{\RExpHaikuPOneWrongDiff}{0.850}
\newcommand{\RExpHaikuPOneWrongHarm}{0}
\newcommand{\RExpHaikuPOneWrongRescue}{51}
\newcommand{\RExpHaikuPThreeCiHi}{0.280}
\newcommand{\RExpHaikuPThreeCiLo}{0.221}
\newcommand{\RExpHaikuPThreeCleanDiff}{0.000}
\newcommand{\RExpHaikuPThreeCleanHarm}{0}
\newcommand{\RExpHaikuPThreeCleanRescue}{0}
\newcommand{\RExpHaikuPThreeDelayedDiff}{$-$0.133}
\newcommand{\RExpHaikuPThreeDelayedHarm}{8}
\newcommand{\RExpHaikuPThreeDelayedHarmTotal}{8}
\newcommand{\RExpHaikuPThreeDelayedRescue}{0}
\newcommand{\RExpHaikuPThreeDiff}{0.250}
\newcommand{\RExpHaikuPThreeDropDiff}{0.933}
\newcommand{\RExpHaikuPThreeDropHarm}{0}
\newcommand{\RExpHaikuPThreeDropRescue}{56}
\newcommand{\RExpHaikuPThreeHarm}{8}
\newcommand{\RExpHaikuPThreeLoudDiff}{0.000}
\newcommand{\RExpHaikuPThreeLoudHarm}{0}
\newcommand{\RExpHaikuPThreeLoudRescue}{0}
\newcommand{\RExpHaikuPThreeNPairs}{380}
\newcommand{\RExpHaikuPThreeP}{$<$0.001}
\newcommand{\RExpHaikuPThreePartialDiff}{0.050}
\newcommand{\RExpHaikuPThreePartialHarm}{0}
\newcommand{\RExpHaikuPThreePartialRescue}{1}
\newcommand{\RExpHaikuPThreeReadDiff}{0.183}
\newcommand{\RExpHaikuPThreeReadHarm}{0}
\newcommand{\RExpHaikuPThreeReadRescue}{11}
\newcommand{\RExpHaikuPThreeRescue}{102}
\newcommand{\RExpHaikuPThreeWrongDiff}{0.567}
\newcommand{\RExpHaikuPThreeWrongHarm}{0}
\newcommand{\RExpHaikuPThreeWrongRescue}{34}
\newcommand{\RExpHaikuPTwoCiHi}{0.283}
\newcommand{\RExpHaikuPTwoCiLo}{0.100}
\newcommand{\RExpHaikuPTwoDiff}{0.183}
\newcommand{\RExpHaikuPTwoHarm}{0}
\newcommand{\RExpHaikuPTwoNPairs}{60}
\newcommand{\RExpHaikuPTwoP}{$<$0.001}
\newcommand{\RExpHaikuPTwoReadDiff}{0.183}
\newcommand{\RExpHaikuPTwoReadHarm}{0}
\newcommand{\RExpHaikuPTwoReadRescue}{11}
\newcommand{\RExpHaikuPTwoRescue}{11}
\newcommand{\RExpHaikuStatusAbstained}{216}
\newcommand{\RExpHaikuStatusBudgetExhausted}{1}
\newcommand{\RExpHaikuStatusDone}{3{,}186}
\newcommand{\RExpHaikuStatusParseFailure}{17}
\newcommand{\RExpQwenEightAppliedAlwaysDelayed}{0.992}
\newcommand{\RExpQwenEightAppliedAlwaysDrop}{0.992}
\newcommand{\RExpQwenEightAppliedAlwaysLoud}{0.992}
\newcommand{\RExpQwenEightAppliedAlwaysPartial}{1.000}
\newcommand{\RExpQwenEightAppliedAlwaysRead}{0.858}
\newcommand{\RExpQwenEightAppliedAlwaysWrong}{0.992}
\newcommand{\RExpQwenEightAppliedBothDelayed}{0.950}
\newcommand{\RExpQwenEightAppliedBothDrop}{0.933}
\newcommand{\RExpQwenEightAppliedBothLoud}{0.950}
\newcommand{\RExpQwenEightAppliedBothPartial}{1.000}
\newcommand{\RExpQwenEightAppliedBothRead}{0.592}
\newcommand{\RExpQwenEightAppliedBothWrong}{0.958}
\newcommand{\RExpQwenEightAppliedContractAbstainDelayed}{0.950}
\newcommand{\RExpQwenEightAppliedContractAbstainDrop}{0.933}
\newcommand{\RExpQwenEightAppliedContractAbstainLoud}{0.950}
\newcommand{\RExpQwenEightAppliedContractAbstainPartial}{1.000}
\newcommand{\RExpQwenEightAppliedContractAbstainRead}{0.592}
\newcommand{\RExpQwenEightAppliedContractAbstainWrong}{0.958}
\newcommand{\RExpQwenEightAppliedContractDelayed}{0.950}
\newcommand{\RExpQwenEightAppliedContractDrop}{0.933}
\newcommand{\RExpQwenEightAppliedContractLoud}{0.950}
\newcommand{\RExpQwenEightAppliedContractPartial}{1.000}
\newcommand{\RExpQwenEightAppliedContractRead}{0.592}
\newcommand{\RExpQwenEightAppliedContractWrong}{0.958}
\newcommand{\RExpQwenEightAppliedNeverDelayed}{0.950}
\newcommand{\RExpQwenEightAppliedNeverDrop}{0.933}
\newcommand{\RExpQwenEightAppliedNeverLoud}{0.950}
\newcommand{\RExpQwenEightAppliedNeverPartial}{1.000}
\newcommand{\RExpQwenEightAppliedNeverRead}{0.592}
\newcommand{\RExpQwenEightAppliedNeverWrong}{0.958}
\newcommand{\RExpQwenEightAppliedOnErrorDelayed}{0.967}
\newcommand{\RExpQwenEightAppliedOnErrorDrop}{0.958}
\newcommand{\RExpQwenEightAppliedOnErrorLoud}{0.975}
\newcommand{\RExpQwenEightAppliedOnErrorPartial}{1.000}
\newcommand{\RExpQwenEightAppliedOnErrorRead}{0.600}
\newcommand{\RExpQwenEightAppliedOnErrorWrong}{0.975}
\newcommand{\RExpQwenEightAppliedOnUncDelayed}{0.950}
\newcommand{\RExpQwenEightAppliedOnUncDrop}{0.942}
\newcommand{\RExpQwenEightAppliedOnUncLoud}{0.958}
\newcommand{\RExpQwenEightAppliedOnUncPartial}{1.000}
\newcommand{\RExpQwenEightAppliedOnUncRead}{0.608}
\newcommand{\RExpQwenEightAppliedOnUncWrong}{0.958}
\newcommand{\RExpQwenEightAppliedReadbackAbstainDelayed}{0.950}
\newcommand{\RExpQwenEightAppliedReadbackAbstainDrop}{0.933}
\newcommand{\RExpQwenEightAppliedReadbackAbstainLoud}{0.950}
\newcommand{\RExpQwenEightAppliedReadbackAbstainPartial}{1.000}
\newcommand{\RExpQwenEightAppliedReadbackAbstainRead}{0.592}
\newcommand{\RExpQwenEightAppliedReadbackAbstainWrong}{0.958}
\newcommand{\RExpQwenEightAppliedReadbackDelayed}{0.950}
\newcommand{\RExpQwenEightAppliedReadbackDrop}{0.933}
\newcommand{\RExpQwenEightAppliedReadbackLoud}{0.950}
\newcommand{\RExpQwenEightAppliedReadbackPartial}{1.000}
\newcommand{\RExpQwenEightAppliedReadbackRead}{0.592}
\newcommand{\RExpQwenEightAppliedReadbackWrong}{0.958}
\newcommand{\RExpQwenEightByKindContractImplausibleFlagged}{33}
\newcommand{\RExpQwenEightByKindContractImplausibleN}{33}
\newcommand{\RExpQwenEightByKindContractImplausibleSuccess}{15}
\newcommand{\RExpQwenEightByKindContractPlausibleFlagged}{0}
\newcommand{\RExpQwenEightByKindContractPlausibleN}{38}
\newcommand{\RExpQwenEightByKindContractPlausibleSuccess}{6}
\newcommand{\RExpQwenEightByKindNeverImplausibleFlagged}{0}
\newcommand{\RExpQwenEightByKindNeverImplausibleN}{33}
\newcommand{\RExpQwenEightByKindNeverImplausibleSuccess}{8}
\newcommand{\RExpQwenEightByKindNeverPlausibleFlagged}{0}
\newcommand{\RExpQwenEightByKindNeverPlausibleN}{38}
\newcommand{\RExpQwenEightByKindNeverPlausibleSuccess}{6}
\newcommand{\RExpQwenEightByKindReadbackImplausibleFlagged}{0}
\newcommand{\RExpQwenEightByKindReadbackImplausibleN}{33}
\newcommand{\RExpQwenEightByKindReadbackImplausibleSuccess}{8}
\newcommand{\RExpQwenEightByKindReadbackPlausibleFlagged}{0}
\newcommand{\RExpQwenEightByKindReadbackPlausibleN}{38}
\newcommand{\RExpQwenEightByKindReadbackPlausibleSuccess}{6}
\newcommand{\RExpQwenEightEpisodes}{6{,}840}
\newcommand{\RExpQwenEightFlagAlways}{1.000}
\newcommand{\RExpQwenEightFlagBoth}{0.107}
\newcommand{\RExpQwenEightFlagContract}{0.000}
\newcommand{\RExpQwenEightFlagContractAbstain}{0.000}
\newcommand{\RExpQwenEightFlagNever}{0.000}
\newcommand{\RExpQwenEightFlagOnError}{0.160}
\newcommand{\RExpQwenEightFlagOnUnc}{0.080}
\newcommand{\RExpQwenEightFlagReadback}{0.107}
\newcommand{\RExpQwenEightFlagReadbackAbstain}{0.135}
\newcommand{\RExpQwenEightPOneCiHi}{0.174}
\newcommand{\RExpQwenEightPOneCiLo}{0.097}
\newcommand{\RExpQwenEightPOneDiff}{0.133}
\newcommand{\RExpQwenEightPOneDropDiff}{0.283}
\newcommand{\RExpQwenEightPOneDropHarm}{0}
\newcommand{\RExpQwenEightPOneDropRescue}{34}
\newcommand{\RExpQwenEightPOneHarm}{0}
\newcommand{\RExpQwenEightPOneNPairs}{280}
\newcommand{\RExpQwenEightPOneP}{$<$0.001}
\newcommand{\RExpQwenEightPOnePartialDiff}{0.000}
\newcommand{\RExpQwenEightPOnePartialHarm}{0}
\newcommand{\RExpQwenEightPOnePartialRescue}{0}
\newcommand{\RExpQwenEightPOneRescue}{35}
\newcommand{\RExpQwenEightPOneWrongDiff}{0.008}
\newcommand{\RExpQwenEightPOneWrongHarm}{0}
\newcommand{\RExpQwenEightPOneWrongRescue}{1}
\newcommand{\RExpQwenEightPThreeCiHi}{0.065}
\newcommand{\RExpQwenEightPThreeCiLo}{0.023}
\newcommand{\RExpQwenEightPThreeCleanDiff}{0.000}
\newcommand{\RExpQwenEightPThreeCleanHarm}{0}
\newcommand{\RExpQwenEightPThreeCleanRescue}{0}
\newcommand{\RExpQwenEightPThreeDelayedDiff}{$-$0.083}
\newcommand{\RExpQwenEightPThreeDelayedHarm}{11}
\newcommand{\RExpQwenEightPThreeDelayedHarmTotal}{11}
\newcommand{\RExpQwenEightPThreeDelayedRescue}{1}
\newcommand{\RExpQwenEightPThreeDiff}{0.044}
\newcommand{\RExpQwenEightPThreeDropDiff}{0.283}
\newcommand{\RExpQwenEightPThreeDropHarm}{0}
\newcommand{\RExpQwenEightPThreeDropRescue}{34}
\newcommand{\RExpQwenEightPThreeHarm}{12}
\newcommand{\RExpQwenEightPThreeLoudDiff}{0.000}
\newcommand{\RExpQwenEightPThreeLoudHarm}{0}
\newcommand{\RExpQwenEightPThreeLoudRescue}{0}
\newcommand{\RExpQwenEightPThreeNPairs}{760}
\newcommand{\RExpQwenEightPThreeP}{$<$0.001}
\newcommand{\RExpQwenEightPThreePartialDiff}{0.000}
\newcommand{\RExpQwenEightPThreePartialHarm}{0}
\newcommand{\RExpQwenEightPThreePartialRescue}{0}
\newcommand{\RExpQwenEightPThreeReadDiff}{0.058}
\newcommand{\RExpQwenEightPThreeReadHarm}{1}
\newcommand{\RExpQwenEightPThreeReadRescue}{8}
\newcommand{\RExpQwenEightPThreeRescue}{44}
\newcommand{\RExpQwenEightPThreeWrongDiff}{0.008}
\newcommand{\RExpQwenEightPThreeWrongHarm}{0}
\newcommand{\RExpQwenEightPThreeWrongRescue}{1}
\newcommand{\RExpQwenEightPTwoCiHi}{0.117}
\newcommand{\RExpQwenEightPTwoCiLo}{0.008}
\newcommand{\RExpQwenEightPTwoDiff}{0.058}
\newcommand{\RExpQwenEightPTwoHarm}{1}
\newcommand{\RExpQwenEightPTwoNPairs}{120}
\newcommand{\RExpQwenEightPTwoP}{0.064}
\newcommand{\RExpQwenEightPTwoReadDiff}{0.058}
\newcommand{\RExpQwenEightPTwoReadHarm}{1}
\newcommand{\RExpQwenEightPTwoReadRescue}{8}
\newcommand{\RExpQwenEightPTwoRescue}{8}
\newcommand{\RExpQwenEightStatusAbstained}{411}
\newcommand{\RExpQwenEightStatusBudgetExhausted}{208}
\newcommand{\RExpQwenEightStatusDone}{5{,}932}
\newcommand{\RExpQwenEightStatusParseFailure}{289}
\newcommand{\RExpQwenFourAppliedAlwaysDelayed}{0.975}
\newcommand{\RExpQwenFourAppliedAlwaysDrop}{0.983}
\newcommand{\RExpQwenFourAppliedAlwaysLoud}{0.975}
\newcommand{\RExpQwenFourAppliedAlwaysPartial}{1.000}
\newcommand{\RExpQwenFourAppliedAlwaysRead}{0.625}
\newcommand{\RExpQwenFourAppliedAlwaysWrong}{0.958}
\newcommand{\RExpQwenFourAppliedBothDelayed}{0.958}
\newcommand{\RExpQwenFourAppliedBothDrop}{0.958}
\newcommand{\RExpQwenFourAppliedBothLoud}{0.958}
\newcommand{\RExpQwenFourAppliedBothPartial}{1.000}
\newcommand{\RExpQwenFourAppliedBothRead}{0.592}
\newcommand{\RExpQwenFourAppliedBothWrong}{0.892}
\newcommand{\RExpQwenFourAppliedContractAbstainDelayed}{0.958}
\newcommand{\RExpQwenFourAppliedContractAbstainDrop}{0.958}
\newcommand{\RExpQwenFourAppliedContractAbstainLoud}{0.958}
\newcommand{\RExpQwenFourAppliedContractAbstainPartial}{1.000}
\newcommand{\RExpQwenFourAppliedContractAbstainRead}{0.592}
\newcommand{\RExpQwenFourAppliedContractAbstainWrong}{0.892}
\newcommand{\RExpQwenFourAppliedContractDelayed}{0.958}
\newcommand{\RExpQwenFourAppliedContractDrop}{0.958}
\newcommand{\RExpQwenFourAppliedContractLoud}{0.958}
\newcommand{\RExpQwenFourAppliedContractPartial}{1.000}
\newcommand{\RExpQwenFourAppliedContractRead}{0.592}
\newcommand{\RExpQwenFourAppliedContractWrong}{0.892}
\newcommand{\RExpQwenFourAppliedNeverDelayed}{0.958}
\newcommand{\RExpQwenFourAppliedNeverDrop}{0.958}
\newcommand{\RExpQwenFourAppliedNeverLoud}{0.958}
\newcommand{\RExpQwenFourAppliedNeverPartial}{1.000}
\newcommand{\RExpQwenFourAppliedNeverRead}{0.592}
\newcommand{\RExpQwenFourAppliedNeverWrong}{0.892}
\newcommand{\RExpQwenFourAppliedOnErrorDelayed}{0.958}
\newcommand{\RExpQwenFourAppliedOnErrorDrop}{0.958}
\newcommand{\RExpQwenFourAppliedOnErrorLoud}{0.967}
\newcommand{\RExpQwenFourAppliedOnErrorPartial}{1.000}
\newcommand{\RExpQwenFourAppliedOnErrorRead}{0.592}
\newcommand{\RExpQwenFourAppliedOnErrorWrong}{0.892}
\newcommand{\RExpQwenFourAppliedOnUncDelayed}{0.967}
\newcommand{\RExpQwenFourAppliedOnUncDrop}{0.967}
\newcommand{\RExpQwenFourAppliedOnUncLoud}{0.967}
\newcommand{\RExpQwenFourAppliedOnUncPartial}{1.000}
\newcommand{\RExpQwenFourAppliedOnUncRead}{0.600}
\newcommand{\RExpQwenFourAppliedOnUncWrong}{0.908}
\newcommand{\RExpQwenFourAppliedReadbackAbstainDelayed}{0.958}
\newcommand{\RExpQwenFourAppliedReadbackAbstainDrop}{0.958}
\newcommand{\RExpQwenFourAppliedReadbackAbstainLoud}{0.958}
\newcommand{\RExpQwenFourAppliedReadbackAbstainPartial}{1.000}
\newcommand{\RExpQwenFourAppliedReadbackAbstainRead}{0.592}
\newcommand{\RExpQwenFourAppliedReadbackAbstainWrong}{0.892}
\newcommand{\RExpQwenFourAppliedReadbackDelayed}{0.958}
\newcommand{\RExpQwenFourAppliedReadbackDrop}{0.958}
\newcommand{\RExpQwenFourAppliedReadbackLoud}{0.958}
\newcommand{\RExpQwenFourAppliedReadbackPartial}{1.000}
\newcommand{\RExpQwenFourAppliedReadbackRead}{0.592}
\newcommand{\RExpQwenFourAppliedReadbackWrong}{0.892}
\newcommand{\RExpQwenFourByKindContractImplausibleFlagged}{30}
\newcommand{\RExpQwenFourByKindContractImplausibleN}{30}
\newcommand{\RExpQwenFourByKindContractImplausibleSuccess}{6}
\newcommand{\RExpQwenFourByKindContractPlausibleFlagged}{0}
\newcommand{\RExpQwenFourByKindContractPlausibleN}{41}
\newcommand{\RExpQwenFourByKindContractPlausibleSuccess}{13}
\newcommand{\RExpQwenFourByKindNeverImplausibleFlagged}{0}
\newcommand{\RExpQwenFourByKindNeverImplausibleN}{30}
\newcommand{\RExpQwenFourByKindNeverImplausibleSuccess}{8}
\newcommand{\RExpQwenFourByKindNeverPlausibleFlagged}{0}
\newcommand{\RExpQwenFourByKindNeverPlausibleN}{41}
\newcommand{\RExpQwenFourByKindNeverPlausibleSuccess}{13}
\newcommand{\RExpQwenFourByKindReadbackImplausibleFlagged}{0}
\newcommand{\RExpQwenFourByKindReadbackImplausibleN}{30}
\newcommand{\RExpQwenFourByKindReadbackImplausibleSuccess}{8}
\newcommand{\RExpQwenFourByKindReadbackPlausibleFlagged}{0}
\newcommand{\RExpQwenFourByKindReadbackPlausibleN}{41}
\newcommand{\RExpQwenFourByKindReadbackPlausibleSuccess}{13}
\newcommand{\RExpQwenFourEpisodes}{6{,}840}
\newcommand{\RExpQwenFourFlagAlways}{1.000}
\newcommand{\RExpQwenFourFlagBoth}{0.114}
\newcommand{\RExpQwenFourFlagContract}{0.000}
\newcommand{\RExpQwenFourFlagContractAbstain}{0.000}
\newcommand{\RExpQwenFourFlagNever}{0.000}
\newcommand{\RExpQwenFourFlagOnError}{0.248}
\newcommand{\RExpQwenFourFlagOnUnc}{0.194}
\newcommand{\RExpQwenFourFlagReadback}{0.114}
\newcommand{\RExpQwenFourFlagReadbackAbstain}{0.131}
\newcommand{\RExpQwenFourPOneCiHi}{0.013}
\newcommand{\RExpQwenFourPOneCiLo}{0.000}
\newcommand{\RExpQwenFourPOneDiff}{0.004}
\newcommand{\RExpQwenFourPOneDropDiff}{0.008}
\newcommand{\RExpQwenFourPOneDropHarm}{0}
\newcommand{\RExpQwenFourPOneDropRescue}{1}
\newcommand{\RExpQwenFourPOneHarm}{0}
\newcommand{\RExpQwenFourPOneNPairs}{280}
\newcommand{\RExpQwenFourPOneP}{1.000}
\newcommand{\RExpQwenFourPOnePartialDiff}{0.000}
\newcommand{\RExpQwenFourPOnePartialHarm}{0}
\newcommand{\RExpQwenFourPOnePartialRescue}{0}
\newcommand{\RExpQwenFourPOneRescue}{1}
\newcommand{\RExpQwenFourPOneWrongDiff}{0.000}
\newcommand{\RExpQwenFourPOneWrongHarm}{0}
\newcommand{\RExpQwenFourPOneWrongRescue}{0}
\newcommand{\RExpQwenFourPThreeCiHi}{$-$0.061}
\newcommand{\RExpQwenFourPThreeCiLo}{$-$0.092}
\newcommand{\RExpQwenFourPThreeCleanDiff}{0.000}
\newcommand{\RExpQwenFourPThreeCleanHarm}{0}
\newcommand{\RExpQwenFourPThreeCleanRescue}{0}
\newcommand{\RExpQwenFourPThreeDelayedDiff}{$-$0.483}
\newcommand{\RExpQwenFourPThreeDelayedHarm}{59}
\newcommand{\RExpQwenFourPThreeDelayedHarmTotal}{59}
\newcommand{\RExpQwenFourPThreeDelayedRescue}{1}
\newcommand{\RExpQwenFourPThreeDiff}{$-$0.077}
\newcommand{\RExpQwenFourPThreeDropDiff}{0.008}
\newcommand{\RExpQwenFourPThreeDropHarm}{0}
\newcommand{\RExpQwenFourPThreeDropRescue}{1}
\newcommand{\RExpQwenFourPThreeHarm}{62}
\newcommand{\RExpQwenFourPThreeLoudDiff}{0.000}
\newcommand{\RExpQwenFourPThreeLoudHarm}{0}
\newcommand{\RExpQwenFourPThreeLoudRescue}{0}
\newcommand{\RExpQwenFourPThreeNPairs}{760}
\newcommand{\RExpQwenFourPThreeP}{$<$0.001}
\newcommand{\RExpQwenFourPThreePartialDiff}{0.000}
\newcommand{\RExpQwenFourPThreePartialHarm}{0}
\newcommand{\RExpQwenFourPThreePartialRescue}{0}
\newcommand{\RExpQwenFourPThreeReadDiff}{$-$0.017}
\newcommand{\RExpQwenFourPThreeReadHarm}{3}
\newcommand{\RExpQwenFourPThreeReadRescue}{1}
\newcommand{\RExpQwenFourPThreeRescue}{3}
\newcommand{\RExpQwenFourPThreeWrongDiff}{0.000}
\newcommand{\RExpQwenFourPThreeWrongHarm}{0}
\newcommand{\RExpQwenFourPThreeWrongRescue}{0}
\newcommand{\RExpQwenFourPTwoCiHi}{0.017}
\newcommand{\RExpQwenFourPTwoCiLo}{$-$0.050}
\newcommand{\RExpQwenFourPTwoDiff}{$-$0.017}
\newcommand{\RExpQwenFourPTwoHarm}{3}
\newcommand{\RExpQwenFourPTwoNPairs}{120}
\newcommand{\RExpQwenFourPTwoP}{0.631}
\newcommand{\RExpQwenFourPTwoReadDiff}{$-$0.017}
\newcommand{\RExpQwenFourPTwoReadHarm}{3}
\newcommand{\RExpQwenFourPTwoReadRescue}{1}
\newcommand{\RExpQwenFourPTwoRescue}{1}
\newcommand{\RExpQwenFourStatusAbstained}{404}
\newcommand{\RExpQwenFourStatusBudgetExhausted}{424}
\newcommand{\RExpQwenFourStatusDone}{4{,}117}
\newcommand{\RExpQwenFourStatusParseFailure}{1{,}895}
\newcommand{\RHaikuCleanAlwaysAbstainRate}{0.000}
\newcommand{\RHaikuCleanAlwaysN}{60}
\newcommand{\RHaikuCleanAlwaysSuccess}{0.933}
\newcommand{\RHaikuCleanAlwaysWrongCommit}{0.050}
\newcommand{\RHaikuCleanBothAbstainRate}{0.000}
\newcommand{\RHaikuCleanBothN}{60}
\newcommand{\RHaikuCleanBothSuccess}{1.000}
\newcommand{\RHaikuCleanBothWrongCommit}{0.000}
\newcommand{\RHaikuCleanContractAbstainAbstainRate}{0.000}
\newcommand{\RHaikuCleanContractAbstainN}{60}
\newcommand{\RHaikuCleanContractAbstainRate}{0.000}
\newcommand{\RHaikuCleanContractAbstainSuccess}{1.000}
\newcommand{\RHaikuCleanContractAbstainWrongCommit}{0.000}
\newcommand{\RHaikuCleanContractN}{60}
\newcommand{\RHaikuCleanContractSuccess}{1.000}
\newcommand{\RHaikuCleanContractWrongCommit}{0.000}
\newcommand{\RHaikuCleanNeverAbstainRate}{0.000}
\newcommand{\RHaikuCleanNeverN}{60}
\newcommand{\RHaikuCleanNeverSuccess}{1.000}
\newcommand{\RHaikuCleanNeverWrongCommit}{0.000}
\newcommand{\RHaikuCleanOnErrorAbstainRate}{0.000}
\newcommand{\RHaikuCleanOnErrorN}{60}
\newcommand{\RHaikuCleanOnErrorSuccess}{1.000}
\newcommand{\RHaikuCleanOnErrorWrongCommit}{0.000}
\newcommand{\RHaikuCleanOnUncAbstainRate}{0.000}
\newcommand{\RHaikuCleanOnUncN}{60}
\newcommand{\RHaikuCleanOnUncSuccess}{1.000}
\newcommand{\RHaikuCleanOnUncWrongCommit}{0.000}
\newcommand{\RHaikuCleanReadbackAbstainAbstainRate}{0.000}
\newcommand{\RHaikuCleanReadbackAbstainN}{60}
\newcommand{\RHaikuCleanReadbackAbstainRate}{0.000}
\newcommand{\RHaikuCleanReadbackAbstainSuccess}{1.000}
\newcommand{\RHaikuCleanReadbackAbstainWrongCommit}{0.000}
\newcommand{\RHaikuCleanReadbackN}{60}
\newcommand{\RHaikuCleanReadbackSuccess}{1.000}
\newcommand{\RHaikuCleanReadbackWrongCommit}{0.000}
\newcommand{\RHaikuCostUsd}{\$0.91}
\newcommand{\RHaikuDelayedAlwaysAbstainRate}{0.000}
\newcommand{\RHaikuDelayedAlwaysN}{60}
\newcommand{\RHaikuDelayedAlwaysSuccess}{0.833}
\newcommand{\RHaikuDelayedAlwaysWrongCommit}{0.117}
\newcommand{\RHaikuDelayedBothAbstainRate}{0.000}
\newcommand{\RHaikuDelayedBothN}{60}
\newcommand{\RHaikuDelayedBothSuccess}{0.867}
\newcommand{\RHaikuDelayedBothWrongCommit}{0.133}
\newcommand{\RHaikuDelayedContractAbstainAbstainRate}{0.000}
\newcommand{\RHaikuDelayedContractAbstainN}{60}
\newcommand{\RHaikuDelayedContractAbstainRate}{0.000}
\newcommand{\RHaikuDelayedContractAbstainSuccess}{1.000}
\newcommand{\RHaikuDelayedContractAbstainWrongCommit}{0.000}
\newcommand{\RHaikuDelayedContractN}{60}
\newcommand{\RHaikuDelayedContractSuccess}{1.000}
\newcommand{\RHaikuDelayedContractWrongCommit}{0.000}
\newcommand{\RHaikuDelayedNeverAbstainRate}{0.000}
\newcommand{\RHaikuDelayedNeverN}{60}
\newcommand{\RHaikuDelayedNeverSuccess}{1.000}
\newcommand{\RHaikuDelayedNeverWrongCommit}{0.000}
\newcommand{\RHaikuDelayedOnErrorAbstainRate}{0.000}
\newcommand{\RHaikuDelayedOnErrorN}{60}
\newcommand{\RHaikuDelayedOnErrorSuccess}{1.000}
\newcommand{\RHaikuDelayedOnErrorWrongCommit}{0.000}
\newcommand{\RHaikuDelayedOnUncAbstainRate}{0.000}
\newcommand{\RHaikuDelayedOnUncN}{60}
\newcommand{\RHaikuDelayedOnUncSuccess}{1.000}
\newcommand{\RHaikuDelayedOnUncWrongCommit}{0.000}
\newcommand{\RHaikuDelayedReadbackAbstainAbstainRate}{1.000}
\newcommand{\RHaikuDelayedReadbackAbstainN}{60}
\newcommand{\RHaikuDelayedReadbackAbstainRate}{0.000}
\newcommand{\RHaikuDelayedReadbackAbstainSuccess}{0.000}
\newcommand{\RHaikuDelayedReadbackAbstainWrongCommit}{0.000}
\newcommand{\RHaikuDelayedReadbackN}{60}
\newcommand{\RHaikuDelayedReadbackSuccess}{0.867}
\newcommand{\RHaikuDelayedReadbackWrongCommit}{0.133}
\newcommand{\RHaikuDropAlwaysAbstainRate}{0.000}
\newcommand{\RHaikuDropAlwaysN}{60}
\newcommand{\RHaikuDropAlwaysSuccess}{0.700}
\newcommand{\RHaikuDropAlwaysWrongCommit}{0.283}
\newcommand{\RHaikuDropBothAbstainRate}{0.000}
\newcommand{\RHaikuDropBothN}{60}
\newcommand{\RHaikuDropBothSuccess}{0.933}
\newcommand{\RHaikuDropBothWrongCommit}{0.067}
\newcommand{\RHaikuDropContractAbstainAbstainRate}{0.000}
\newcommand{\RHaikuDropContractAbstainN}{60}
\newcommand{\RHaikuDropContractAbstainRate}{0.000}
\newcommand{\RHaikuDropContractAbstainSuccess}{0.000}
\newcommand{\RHaikuDropContractAbstainWrongCommit}{1.000}
\newcommand{\RHaikuDropContractN}{60}
\newcommand{\RHaikuDropContractSuccess}{0.000}
\newcommand{\RHaikuDropContractWrongCommit}{1.000}
\newcommand{\RHaikuDropNeverAbstainRate}{0.000}
\newcommand{\RHaikuDropNeverN}{60}
\newcommand{\RHaikuDropNeverSuccess}{0.000}
\newcommand{\RHaikuDropNeverWrongCommit}{1.000}
\newcommand{\RHaikuDropOnErrorAbstainRate}{0.000}
\newcommand{\RHaikuDropOnErrorN}{60}
\newcommand{\RHaikuDropOnErrorSuccess}{0.000}
\newcommand{\RHaikuDropOnErrorWrongCommit}{1.000}
\newcommand{\RHaikuDropOnUncAbstainRate}{0.000}
\newcommand{\RHaikuDropOnUncN}{60}
\newcommand{\RHaikuDropOnUncSuccess}{0.050}
\newcommand{\RHaikuDropOnUncWrongCommit}{0.950}
\newcommand{\RHaikuDropReadbackAbstainAbstainRate}{1.000}
\newcommand{\RHaikuDropReadbackAbstainN}{60}
\newcommand{\RHaikuDropReadbackAbstainRate}{0.000}
\newcommand{\RHaikuDropReadbackAbstainSuccess}{0.000}
\newcommand{\RHaikuDropReadbackAbstainWrongCommit}{0.000}
\newcommand{\RHaikuDropReadbackN}{60}
\newcommand{\RHaikuDropReadbackSuccess}{0.933}
\newcommand{\RHaikuDropReadbackWrongCommit}{0.067}
\newcommand{\RHaikuLoudAlwaysAbstainRate}{0.000}
\newcommand{\RHaikuLoudAlwaysN}{60}
\newcommand{\RHaikuLoudAlwaysSuccess}{0.833}
\newcommand{\RHaikuLoudAlwaysWrongCommit}{0.150}
\newcommand{\RHaikuLoudBothAbstainRate}{0.000}
\newcommand{\RHaikuLoudBothN}{60}
\newcommand{\RHaikuLoudBothSuccess}{1.000}
\newcommand{\RHaikuLoudBothWrongCommit}{0.000}
\newcommand{\RHaikuLoudContractAbstainAbstainRate}{0.000}
\newcommand{\RHaikuLoudContractAbstainN}{60}
\newcommand{\RHaikuLoudContractAbstainRate}{0.000}
\newcommand{\RHaikuLoudContractAbstainSuccess}{1.000}
\newcommand{\RHaikuLoudContractAbstainWrongCommit}{0.000}
\newcommand{\RHaikuLoudContractN}{60}
\newcommand{\RHaikuLoudContractSuccess}{1.000}
\newcommand{\RHaikuLoudContractWrongCommit}{0.000}
\newcommand{\RHaikuLoudNeverAbstainRate}{0.000}
\newcommand{\RHaikuLoudNeverN}{60}
\newcommand{\RHaikuLoudNeverSuccess}{1.000}
\newcommand{\RHaikuLoudNeverWrongCommit}{0.000}
\newcommand{\RHaikuLoudOnErrorAbstainRate}{0.000}
\newcommand{\RHaikuLoudOnErrorN}{60}
\newcommand{\RHaikuLoudOnErrorSuccess}{0.917}
\newcommand{\RHaikuLoudOnErrorWrongCommit}{0.083}
\newcommand{\RHaikuLoudOnUncAbstainRate}{0.000}
\newcommand{\RHaikuLoudOnUncN}{60}
\newcommand{\RHaikuLoudOnUncSuccess}{0.983}
\newcommand{\RHaikuLoudOnUncWrongCommit}{0.000}
\newcommand{\RHaikuLoudReadbackAbstainAbstainRate}{0.000}
\newcommand{\RHaikuLoudReadbackAbstainN}{60}
\newcommand{\RHaikuLoudReadbackAbstainRate}{0.000}
\newcommand{\RHaikuLoudReadbackAbstainSuccess}{1.000}
\newcommand{\RHaikuLoudReadbackAbstainWrongCommit}{0.000}
\newcommand{\RHaikuLoudReadbackN}{60}
\newcommand{\RHaikuLoudReadbackSuccess}{1.000}
\newcommand{\RHaikuLoudReadbackWrongCommit}{0.000}
\newcommand{\RHaikuMaxSteps}{15}
\newcommand{\RHaikuModel}{claude-haiku-5-5}
\newcommand{\RHaikuPartialAlwaysAbstainRate}{0.000}
\newcommand{\RHaikuPartialAlwaysN}{20}
\newcommand{\RHaikuPartialAlwaysSuccess}{0.150}
\newcommand{\RHaikuPartialAlwaysWrongCommit}{0.700}
\newcommand{\RHaikuPartialBothAbstainRate}{0.000}
\newcommand{\RHaikuPartialBothN}{20}
\newcommand{\RHaikuPartialBothSuccess}{0.050}
\newcommand{\RHaikuPartialBothWrongCommit}{0.900}
\newcommand{\RHaikuPartialContractAbstainAbstainRate}{0.000}
\newcommand{\RHaikuPartialContractAbstainN}{20}
\newcommand{\RHaikuPartialContractAbstainRate}{0.000}
\newcommand{\RHaikuPartialContractAbstainSuccess}{0.000}
\newcommand{\RHaikuPartialContractAbstainWrongCommit}{1.000}
\newcommand{\RHaikuPartialContractN}{20}
\newcommand{\RHaikuPartialContractSuccess}{0.000}
\newcommand{\RHaikuPartialContractWrongCommit}{1.000}
\newcommand{\RHaikuPartialNeverAbstainRate}{0.000}
\newcommand{\RHaikuPartialNeverN}{20}
\newcommand{\RHaikuPartialNeverSuccess}{0.000}
\newcommand{\RHaikuPartialNeverWrongCommit}{1.000}
\newcommand{\RHaikuPartialOnErrorAbstainRate}{0.000}
\newcommand{\RHaikuPartialOnErrorN}{20}
\newcommand{\RHaikuPartialOnErrorSuccess}{0.000}
\newcommand{\RHaikuPartialOnErrorWrongCommit}{1.000}
\newcommand{\RHaikuPartialOnUncAbstainRate}{0.000}
\newcommand{\RHaikuPartialOnUncN}{20}
\newcommand{\RHaikuPartialOnUncSuccess}{0.000}
\newcommand{\RHaikuPartialOnUncWrongCommit}{1.000}
\newcommand{\RHaikuPartialReadbackAbstainAbstainRate}{1.000}
\newcommand{\RHaikuPartialReadbackAbstainN}{20}
\newcommand{\RHaikuPartialReadbackAbstainRate}{0.000}
\newcommand{\RHaikuPartialReadbackAbstainSuccess}{0.000}
\newcommand{\RHaikuPartialReadbackAbstainWrongCommit}{0.000}
\newcommand{\RHaikuPartialReadbackN}{20}
\newcommand{\RHaikuPartialReadbackSuccess}{0.450}
\newcommand{\RHaikuPartialReadbackWrongCommit}{0.550}
\newcommand{\RHaikuReadAlwaysAbstainRate}{0.000}
\newcommand{\RHaikuReadAlwaysN}{60}
\newcommand{\RHaikuReadAlwaysSuccess}{0.733}
\newcommand{\RHaikuReadAlwaysWrongCommit}{0.217}
\newcommand{\RHaikuReadBothAbstainRate}{0.000}
\newcommand{\RHaikuReadBothN}{60}
\newcommand{\RHaikuReadBothSuccess}{0.783}
\newcommand{\RHaikuReadBothWrongCommit}{0.217}
\newcommand{\RHaikuReadContractAbstainAbstainRate}{0.267}
\newcommand{\RHaikuReadContractAbstainN}{60}
\newcommand{\RHaikuReadContractAbstainRate}{0.000}
\newcommand{\RHaikuReadContractAbstainSuccess}{0.517}
\newcommand{\RHaikuReadContractAbstainWrongCommit}{0.217}
\newcommand{\RHaikuReadContractN}{60}
\newcommand{\RHaikuReadContractSuccess}{0.783}
\newcommand{\RHaikuReadContractWrongCommit}{0.217}
\newcommand{\RHaikuReadNeverAbstainRate}{0.000}
\newcommand{\RHaikuReadNeverN}{60}
\newcommand{\RHaikuReadNeverSuccess}{0.600}
\newcommand{\RHaikuReadNeverWrongCommit}{0.400}
\newcommand{\RHaikuReadOnErrorAbstainRate}{0.000}
\newcommand{\RHaikuReadOnErrorN}{60}
\newcommand{\RHaikuReadOnErrorSuccess}{0.633}
\newcommand{\RHaikuReadOnErrorWrongCommit}{0.367}
\newcommand{\RHaikuReadOnUncAbstainRate}{0.000}
\newcommand{\RHaikuReadOnUncN}{60}
\newcommand{\RHaikuReadOnUncSuccess}{0.583}
\newcommand{\RHaikuReadOnUncWrongCommit}{0.400}
\newcommand{\RHaikuReadReadbackAbstainAbstainRate}{0.000}
\newcommand{\RHaikuReadReadbackAbstainN}{60}
\newcommand{\RHaikuReadReadbackAbstainRate}{0.000}
\newcommand{\RHaikuReadReadbackAbstainSuccess}{0.600}
\newcommand{\RHaikuReadReadbackAbstainWrongCommit}{0.400}
\newcommand{\RHaikuReadReadbackN}{60}
\newcommand{\RHaikuReadReadbackSuccess}{0.600}
\newcommand{\RHaikuReadReadbackWrongCommit}{0.400}
\newcommand{\RHaikuSeedFirst}{100}
\newcommand{\RHaikuThreshold}{0.900}
\newcommand{\RHaikuWrongAlwaysAbstainRate}{0.000}
\newcommand{\RHaikuWrongAlwaysN}{60}
\newcommand{\RHaikuWrongAlwaysSuccess}{0.133}
\newcommand{\RHaikuWrongAlwaysWrongCommit}{0.850}
\newcommand{\RHaikuWrongBothAbstainRate}{0.000}
\newcommand{\RHaikuWrongBothN}{60}
\newcommand{\RHaikuWrongBothSuccess}{0.567}
\newcommand{\RHaikuWrongBothWrongCommit}{0.400}
\newcommand{\RHaikuWrongContractAbstainAbstainRate}{0.000}
\newcommand{\RHaikuWrongContractAbstainN}{60}
\newcommand{\RHaikuWrongContractAbstainRate}{0.000}
\newcommand{\RHaikuWrongContractAbstainSuccess}{0.000}
\newcommand{\RHaikuWrongContractAbstainWrongCommit}{1.000}
\newcommand{\RHaikuWrongContractN}{60}
\newcommand{\RHaikuWrongContractSuccess}{0.000}
\newcommand{\RHaikuWrongContractWrongCommit}{1.000}
\newcommand{\RHaikuWrongNeverAbstainRate}{0.000}
\newcommand{\RHaikuWrongNeverN}{60}
\newcommand{\RHaikuWrongNeverSuccess}{0.000}
\newcommand{\RHaikuWrongNeverWrongCommit}{1.000}
\newcommand{\RHaikuWrongOnErrorAbstainRate}{0.000}
\newcommand{\RHaikuWrongOnErrorN}{60}
\newcommand{\RHaikuWrongOnErrorSuccess}{0.000}
\newcommand{\RHaikuWrongOnErrorWrongCommit}{1.000}
\newcommand{\RHaikuWrongOnUncAbstainRate}{0.000}
\newcommand{\RHaikuWrongOnUncN}{60}
\newcommand{\RHaikuWrongOnUncSuccess}{0.017}
\newcommand{\RHaikuWrongOnUncWrongCommit}{0.983}
\newcommand{\RHaikuWrongReadbackAbstainAbstainRate}{1.000}
\newcommand{\RHaikuWrongReadbackAbstainN}{60}
\newcommand{\RHaikuWrongReadbackAbstainRate}{0.000}
\newcommand{\RHaikuWrongReadbackAbstainSuccess}{0.000}
\newcommand{\RHaikuWrongReadbackAbstainWrongCommit}{0.000}
\newcommand{\RHaikuWrongReadbackN}{60}
\newcommand{\RHaikuWrongReadbackSuccess}{0.850}
\newcommand{\RHaikuWrongReadbackWrongCommit}{0.150}
\newcommand{\RPOneCiHi}{0.174}
\newcommand{\RPOneCiLo}{0.097}
\newcommand{\RPOneDiff}{0.133}
\newcommand{\RPOneDropDiff}{0.283}
\newcommand{\RPOneDropHarm}{0}
\newcommand{\RPOneDropRescue}{34}
\newcommand{\RPOneHarm}{0}
\newcommand{\RPOneHolmP}{$<$0.001}
\newcommand{\RPOneNPairs}{280}
\newcommand{\RPOneP}{$<$0.001}
\newcommand{\RPOnePMcnemar}{$<$0.001}
\newcommand{\RPOnePartialDiff}{0.000}
\newcommand{\RPOnePartialHarm}{0}
\newcommand{\RPOnePartialRescue}{0}
\newcommand{\RPOneRescue}{35}
\newcommand{\RPOneWrongDiff}{0.008}
\newcommand{\RPOneWrongHarm}{0}
\newcommand{\RPOneWrongRescue}{1}
\newcommand{\RPThreeCiHi}{0.065}
\newcommand{\RPThreeCiLo}{0.023}
\newcommand{\RPThreeCleanDiff}{0.000}
\newcommand{\RPThreeCleanHarm}{0}
\newcommand{\RPThreeCleanRescue}{0}
\newcommand{\RPThreeDelayedDiff}{$-$0.083}
\newcommand{\RPThreeDelayedHarm}{11}
\newcommand{\RPThreeDelayedHarmTotal}{11}
\newcommand{\RPThreeDelayedRescue}{1}
\newcommand{\RPThreeDiff}{0.044}
\newcommand{\RPThreeDropDiff}{0.283}
\newcommand{\RPThreeDropHarm}{0}
\newcommand{\RPThreeDropRescue}{34}
\newcommand{\RPThreeHarm}{12}
\newcommand{\RPThreeHolmP}{0.001}
\newcommand{\RPThreeLoudDiff}{0.000}
\newcommand{\RPThreeLoudHarm}{0}
\newcommand{\RPThreeLoudRescue}{0}
\newcommand{\RPThreeNPairs}{760}
\newcommand{\RPThreeP}{$<$0.001}
\newcommand{\RPThreePMcnemar}{$<$0.001}
\newcommand{\RPThreePartialDiff}{0.000}
\newcommand{\RPThreePartialHarm}{0}
\newcommand{\RPThreePartialRescue}{0}
\newcommand{\RPThreeReadDiff}{0.058}
\newcommand{\RPThreeReadHarm}{1}
\newcommand{\RPThreeReadRescue}{8}
\newcommand{\RPThreeRescue}{44}
\newcommand{\RPThreeWrongDiff}{0.008}
\newcommand{\RPThreeWrongHarm}{0}
\newcommand{\RPThreeWrongRescue}{1}
\newcommand{\RPTwoCiHi}{0.117}
\newcommand{\RPTwoCiLo}{0.008}
\newcommand{\RPTwoDiff}{0.058}
\newcommand{\RPTwoHarm}{1}
\newcommand{\RPTwoHolmP}{0.064}
\newcommand{\RPTwoNPairs}{120}
\newcommand{\RPTwoP}{0.064}
\newcommand{\RPTwoPMcnemar}{0.039}
\newcommand{\RPTwoReadDiff}{0.058}
\newcommand{\RPTwoReadHarm}{1}
\newcommand{\RPTwoReadRescue}{8}
\newcommand{\RPTwoRescue}{8}
\newcommand{\RQwenEightCleanAlwaysAbstainRate}{0.000}
\newcommand{\RQwenEightCleanAlwaysN}{120}
\newcommand{\RQwenEightCleanAlwaysSuccess}{0.692}
\newcommand{\RQwenEightCleanAlwaysWrongCommit}{0.100}
\newcommand{\RQwenEightCleanBothAbstainRate}{0.000}
\newcommand{\RQwenEightCleanBothN}{120}
\newcommand{\RQwenEightCleanBothSuccess}{0.767}
\newcommand{\RQwenEightCleanBothWrongCommit}{0.208}
\newcommand{\RQwenEightCleanContractAbstainAbstainRate}{0.000}
\newcommand{\RQwenEightCleanContractAbstainN}{120}
\newcommand{\RQwenEightCleanContractAbstainRate}{0.000}
\newcommand{\RQwenEightCleanContractAbstainSuccess}{0.767}
\newcommand{\RQwenEightCleanContractAbstainWrongCommit}{0.208}
\newcommand{\RQwenEightCleanContractN}{120}
\newcommand{\RQwenEightCleanContractSuccess}{0.767}
\newcommand{\RQwenEightCleanContractWrongCommit}{0.208}
\newcommand{\RQwenEightCleanNeverAbstainRate}{0.000}
\newcommand{\RQwenEightCleanNeverN}{120}
\newcommand{\RQwenEightCleanNeverSuccess}{0.767}
\newcommand{\RQwenEightCleanNeverWrongCommit}{0.208}
\newcommand{\RQwenEightCleanOnErrorAbstainRate}{0.000}
\newcommand{\RQwenEightCleanOnErrorN}{120}
\newcommand{\RQwenEightCleanOnErrorSuccess}{0.775}
\newcommand{\RQwenEightCleanOnErrorWrongCommit}{0.225}
\newcommand{\RQwenEightCleanOnUncAbstainRate}{0.000}
\newcommand{\RQwenEightCleanOnUncN}{120}
\newcommand{\RQwenEightCleanOnUncSuccess}{0.775}
\newcommand{\RQwenEightCleanOnUncWrongCommit}{0.200}
\newcommand{\RQwenEightCleanReadbackAbstainAbstainRate}{0.000}
\newcommand{\RQwenEightCleanReadbackAbstainN}{120}
\newcommand{\RQwenEightCleanReadbackAbstainRate}{0.000}
\newcommand{\RQwenEightCleanReadbackAbstainSuccess}{0.767}
\newcommand{\RQwenEightCleanReadbackAbstainWrongCommit}{0.208}
\newcommand{\RQwenEightCleanReadbackN}{120}
\newcommand{\RQwenEightCleanReadbackSuccess}{0.767}
\newcommand{\RQwenEightCleanReadbackWrongCommit}{0.208}
\newcommand{\RQwenEightDelayedAlwaysAbstainRate}{0.000}
\newcommand{\RQwenEightDelayedAlwaysN}{120}
\newcommand{\RQwenEightDelayedAlwaysSuccess}{0.667}
\newcommand{\RQwenEightDelayedAlwaysWrongCommit}{0.108}
\newcommand{\RQwenEightDelayedBothAbstainRate}{0.000}
\newcommand{\RQwenEightDelayedBothN}{120}
\newcommand{\RQwenEightDelayedBothSuccess}{0.683}
\newcommand{\RQwenEightDelayedBothWrongCommit}{0.183}
\newcommand{\RQwenEightDelayedContractAbstainAbstainRate}{0.000}
\newcommand{\RQwenEightDelayedContractAbstainN}{120}
\newcommand{\RQwenEightDelayedContractAbstainRate}{0.000}
\newcommand{\RQwenEightDelayedContractAbstainSuccess}{0.767}
\newcommand{\RQwenEightDelayedContractAbstainWrongCommit}{0.208}
\newcommand{\RQwenEightDelayedContractN}{120}
\newcommand{\RQwenEightDelayedContractSuccess}{0.767}
\newcommand{\RQwenEightDelayedContractWrongCommit}{0.208}
\newcommand{\RQwenEightDelayedNeverAbstainRate}{0.000}
\newcommand{\RQwenEightDelayedNeverN}{120}
\newcommand{\RQwenEightDelayedNeverSuccess}{0.767}
\newcommand{\RQwenEightDelayedNeverWrongCommit}{0.208}
\newcommand{\RQwenEightDelayedOnErrorAbstainRate}{0.000}
\newcommand{\RQwenEightDelayedOnErrorN}{120}
\newcommand{\RQwenEightDelayedOnErrorSuccess}{0.775}
\newcommand{\RQwenEightDelayedOnErrorWrongCommit}{0.225}
\newcommand{\RQwenEightDelayedOnUncAbstainRate}{0.000}
\newcommand{\RQwenEightDelayedOnUncN}{120}
\newcommand{\RQwenEightDelayedOnUncSuccess}{0.775}
\newcommand{\RQwenEightDelayedOnUncWrongCommit}{0.192}
\newcommand{\RQwenEightDelayedReadbackAbstainAbstainRate}{0.950}
\newcommand{\RQwenEightDelayedReadbackAbstainN}{120}
\newcommand{\RQwenEightDelayedReadbackAbstainRate}{0.000}
\newcommand{\RQwenEightDelayedReadbackAbstainSuccess}{0.000}
\newcommand{\RQwenEightDelayedReadbackAbstainWrongCommit}{0.033}
\newcommand{\RQwenEightDelayedReadbackN}{120}
\newcommand{\RQwenEightDelayedReadbackSuccess}{0.683}
\newcommand{\RQwenEightDelayedReadbackWrongCommit}{0.183}
\newcommand{\RQwenEightDropAlwaysAbstainRate}{0.000}
\newcommand{\RQwenEightDropAlwaysN}{120}
\newcommand{\RQwenEightDropAlwaysSuccess}{0.133}
\newcommand{\RQwenEightDropAlwaysWrongCommit}{0.617}
\newcommand{\RQwenEightDropBothAbstainRate}{0.000}
\newcommand{\RQwenEightDropBothN}{120}
\newcommand{\RQwenEightDropBothSuccess}{0.300}
\newcommand{\RQwenEightDropBothWrongCommit}{0.567}
\newcommand{\RQwenEightDropContractAbstainAbstainRate}{0.000}
\newcommand{\RQwenEightDropContractAbstainN}{120}
\newcommand{\RQwenEightDropContractAbstainRate}{0.000}
\newcommand{\RQwenEightDropContractAbstainSuccess}{0.017}
\newcommand{\RQwenEightDropContractAbstainWrongCommit}{0.958}
\newcommand{\RQwenEightDropContractN}{120}
\newcommand{\RQwenEightDropContractSuccess}{0.017}
\newcommand{\RQwenEightDropContractWrongCommit}{0.958}
\newcommand{\RQwenEightDropNeverAbstainRate}{0.000}
\newcommand{\RQwenEightDropNeverN}{120}
\newcommand{\RQwenEightDropNeverSuccess}{0.017}
\newcommand{\RQwenEightDropNeverWrongCommit}{0.958}
\newcommand{\RQwenEightDropOnErrorAbstainRate}{0.000}
\newcommand{\RQwenEightDropOnErrorN}{120}
\newcommand{\RQwenEightDropOnErrorSuccess}{0.008}
\newcommand{\RQwenEightDropOnErrorWrongCommit}{0.992}
\newcommand{\RQwenEightDropOnUncAbstainRate}{0.000}
\newcommand{\RQwenEightDropOnUncN}{120}
\newcommand{\RQwenEightDropOnUncSuccess}{0.017}
\newcommand{\RQwenEightDropOnUncWrongCommit}{0.958}
\newcommand{\RQwenEightDropReadbackAbstainAbstainRate}{0.917}
\newcommand{\RQwenEightDropReadbackAbstainN}{120}
\newcommand{\RQwenEightDropReadbackAbstainRate}{0.000}
\newcommand{\RQwenEightDropReadbackAbstainSuccess}{0.017}
\newcommand{\RQwenEightDropReadbackAbstainWrongCommit}{0.050}
\newcommand{\RQwenEightDropReadbackN}{120}
\newcommand{\RQwenEightDropReadbackSuccess}{0.300}
\newcommand{\RQwenEightDropReadbackWrongCommit}{0.567}
\newcommand{\RQwenEightLoudAlwaysAbstainRate}{0.000}
\newcommand{\RQwenEightLoudAlwaysN}{120}
\newcommand{\RQwenEightLoudAlwaysSuccess}{0.683}
\newcommand{\RQwenEightLoudAlwaysWrongCommit}{0.117}
\newcommand{\RQwenEightLoudBothAbstainRate}{0.000}
\newcommand{\RQwenEightLoudBothN}{120}
\newcommand{\RQwenEightLoudBothSuccess}{0.742}
\newcommand{\RQwenEightLoudBothWrongCommit}{0.200}
\newcommand{\RQwenEightLoudContractAbstainAbstainRate}{0.000}
\newcommand{\RQwenEightLoudContractAbstainN}{120}
\newcommand{\RQwenEightLoudContractAbstainRate}{0.000}
\newcommand{\RQwenEightLoudContractAbstainSuccess}{0.742}
\newcommand{\RQwenEightLoudContractAbstainWrongCommit}{0.200}
\newcommand{\RQwenEightLoudContractN}{120}
\newcommand{\RQwenEightLoudContractSuccess}{0.742}
\newcommand{\RQwenEightLoudContractWrongCommit}{0.200}
\newcommand{\RQwenEightLoudNeverAbstainRate}{0.000}
\newcommand{\RQwenEightLoudNeverN}{120}
\newcommand{\RQwenEightLoudNeverSuccess}{0.742}
\newcommand{\RQwenEightLoudNeverWrongCommit}{0.200}
\newcommand{\RQwenEightLoudOnErrorAbstainRate}{0.000}
\newcommand{\RQwenEightLoudOnErrorN}{120}
\newcommand{\RQwenEightLoudOnErrorSuccess}{0.783}
\newcommand{\RQwenEightLoudOnErrorWrongCommit}{0.175}
\newcommand{\RQwenEightLoudOnUncAbstainRate}{0.000}
\newcommand{\RQwenEightLoudOnUncN}{120}
\newcommand{\RQwenEightLoudOnUncSuccess}{0.600}
\newcommand{\RQwenEightLoudOnUncWrongCommit}{0.192}
\newcommand{\RQwenEightLoudReadbackAbstainAbstainRate}{0.000}
\newcommand{\RQwenEightLoudReadbackAbstainN}{120}
\newcommand{\RQwenEightLoudReadbackAbstainRate}{0.000}
\newcommand{\RQwenEightLoudReadbackAbstainSuccess}{0.742}
\newcommand{\RQwenEightLoudReadbackAbstainWrongCommit}{0.200}
\newcommand{\RQwenEightLoudReadbackN}{120}
\newcommand{\RQwenEightLoudReadbackSuccess}{0.742}
\newcommand{\RQwenEightLoudReadbackWrongCommit}{0.200}
\newcommand{\RQwenEightMaxSteps}{15}
\newcommand{\RQwenEightModel}{qwen3:8b-q4\_K\_M}
\newcommand{\RQwenEightPartialAlwaysAbstainRate}{0.000}
\newcommand{\RQwenEightPartialAlwaysN}{40}
\newcommand{\RQwenEightPartialAlwaysSuccess}{0.000}
\newcommand{\RQwenEightPartialAlwaysWrongCommit}{0.700}
\newcommand{\RQwenEightPartialBothAbstainRate}{0.000}
\newcommand{\RQwenEightPartialBothN}{40}
\newcommand{\RQwenEightPartialBothSuccess}{0.000}
\newcommand{\RQwenEightPartialBothWrongCommit}{0.975}
\newcommand{\RQwenEightPartialContractAbstainAbstainRate}{0.000}
\newcommand{\RQwenEightPartialContractAbstainN}{40}
\newcommand{\RQwenEightPartialContractAbstainRate}{0.000}
\newcommand{\RQwenEightPartialContractAbstainSuccess}{0.000}
\newcommand{\RQwenEightPartialContractAbstainWrongCommit}{1.000}
\newcommand{\RQwenEightPartialContractN}{40}
\newcommand{\RQwenEightPartialContractSuccess}{0.000}
\newcommand{\RQwenEightPartialContractWrongCommit}{1.000}
\newcommand{\RQwenEightPartialNeverAbstainRate}{0.000}
\newcommand{\RQwenEightPartialNeverN}{40}
\newcommand{\RQwenEightPartialNeverSuccess}{0.000}
\newcommand{\RQwenEightPartialNeverWrongCommit}{1.000}
\newcommand{\RQwenEightPartialOnErrorAbstainRate}{0.000}
\newcommand{\RQwenEightPartialOnErrorN}{40}
\newcommand{\RQwenEightPartialOnErrorSuccess}{0.000}
\newcommand{\RQwenEightPartialOnErrorWrongCommit}{1.000}
\newcommand{\RQwenEightPartialOnUncAbstainRate}{0.000}
\newcommand{\RQwenEightPartialOnUncN}{40}
\newcommand{\RQwenEightPartialOnUncSuccess}{0.000}
\newcommand{\RQwenEightPartialOnUncWrongCommit}{1.000}
\newcommand{\RQwenEightPartialReadbackAbstainAbstainRate}{1.000}
\newcommand{\RQwenEightPartialReadbackAbstainN}{40}
\newcommand{\RQwenEightPartialReadbackAbstainRate}{0.000}
\newcommand{\RQwenEightPartialReadbackAbstainSuccess}{0.000}
\newcommand{\RQwenEightPartialReadbackAbstainWrongCommit}{0.000}
\newcommand{\RQwenEightPartialReadbackN}{40}
\newcommand{\RQwenEightPartialReadbackSuccess}{0.000}
\newcommand{\RQwenEightPartialReadbackWrongCommit}{0.975}
\newcommand{\RQwenEightReadAlwaysAbstainRate}{0.000}
\newcommand{\RQwenEightReadAlwaysN}{120}
\newcommand{\RQwenEightReadAlwaysSuccess}{0.325}
\newcommand{\RQwenEightReadAlwaysWrongCommit}{0.400}
\newcommand{\RQwenEightReadBothAbstainRate}{0.000}
\newcommand{\RQwenEightReadBothN}{120}
\newcommand{\RQwenEightReadBothSuccess}{0.508}
\newcommand{\RQwenEightReadBothWrongCommit}{0.417}
\newcommand{\RQwenEightReadContractAbstainAbstainRate}{0.275}
\newcommand{\RQwenEightReadContractAbstainN}{120}
\newcommand{\RQwenEightReadContractAbstainRate}{0.000}
\newcommand{\RQwenEightReadContractAbstainSuccess}{0.383}
\newcommand{\RQwenEightReadContractAbstainWrongCommit}{0.317}
\newcommand{\RQwenEightReadContractN}{120}
\newcommand{\RQwenEightReadContractSuccess}{0.508}
\newcommand{\RQwenEightReadContractWrongCommit}{0.417}
\newcommand{\RQwenEightReadNeverAbstainRate}{0.000}
\newcommand{\RQwenEightReadNeverN}{120}
\newcommand{\RQwenEightReadNeverSuccess}{0.450}
\newcommand{\RQwenEightReadNeverWrongCommit}{0.500}
\newcommand{\RQwenEightReadOnErrorAbstainRate}{0.000}
\newcommand{\RQwenEightReadOnErrorN}{120}
\newcommand{\RQwenEightReadOnErrorSuccess}{0.450}
\newcommand{\RQwenEightReadOnErrorWrongCommit}{0.525}
\newcommand{\RQwenEightReadOnUncAbstainRate}{0.000}
\newcommand{\RQwenEightReadOnUncN}{120}
\newcommand{\RQwenEightReadOnUncSuccess}{0.458}
\newcommand{\RQwenEightReadOnUncWrongCommit}{0.508}
\newcommand{\RQwenEightReadReadbackAbstainAbstainRate}{0.000}
\newcommand{\RQwenEightReadReadbackAbstainN}{120}
\newcommand{\RQwenEightReadReadbackAbstainRate}{0.000}
\newcommand{\RQwenEightReadReadbackAbstainSuccess}{0.450}
\newcommand{\RQwenEightReadReadbackAbstainWrongCommit}{0.500}
\newcommand{\RQwenEightReadReadbackN}{120}
\newcommand{\RQwenEightReadReadbackSuccess}{0.450}
\newcommand{\RQwenEightReadReadbackWrongCommit}{0.500}
\newcommand{\RQwenEightSeedFirst}{100}
\newcommand{\RQwenEightSeedSecond}{101}
\newcommand{\RQwenEightThreshold}{0.900}
\newcommand{\RQwenEightWrongAlwaysAbstainRate}{0.000}
\newcommand{\RQwenEightWrongAlwaysN}{120}
\newcommand{\RQwenEightWrongAlwaysSuccess}{0.017}
\newcommand{\RQwenEightWrongAlwaysWrongCommit}{0.750}
\newcommand{\RQwenEightWrongBothAbstainRate}{0.000}
\newcommand{\RQwenEightWrongBothN}{120}
\newcommand{\RQwenEightWrongBothSuccess}{0.008}
\newcommand{\RQwenEightWrongBothWrongCommit}{0.708}
\newcommand{\RQwenEightWrongContractAbstainAbstainRate}{0.008}
\newcommand{\RQwenEightWrongContractAbstainN}{120}
\newcommand{\RQwenEightWrongContractAbstainRate}{0.000}
\newcommand{\RQwenEightWrongContractAbstainSuccess}{0.000}
\newcommand{\RQwenEightWrongContractAbstainWrongCommit}{0.967}
\newcommand{\RQwenEightWrongContractN}{120}
\newcommand{\RQwenEightWrongContractSuccess}{0.000}
\newcommand{\RQwenEightWrongContractWrongCommit}{0.975}
\newcommand{\RQwenEightWrongNeverAbstainRate}{0.000}
\newcommand{\RQwenEightWrongNeverN}{120}
\newcommand{\RQwenEightWrongNeverSuccess}{0.000}
\newcommand{\RQwenEightWrongNeverWrongCommit}{0.975}
\newcommand{\RQwenEightWrongOnErrorAbstainRate}{0.000}
\newcommand{\RQwenEightWrongOnErrorN}{120}
\newcommand{\RQwenEightWrongOnErrorSuccess}{0.000}
\newcommand{\RQwenEightWrongOnErrorWrongCommit}{1.000}
\newcommand{\RQwenEightWrongOnUncAbstainRate}{0.000}
\newcommand{\RQwenEightWrongOnUncN}{120}
\newcommand{\RQwenEightWrongOnUncSuccess}{0.000}
\newcommand{\RQwenEightWrongOnUncWrongCommit}{0.975}
\newcommand{\RQwenEightWrongReadbackAbstainAbstainRate}{0.942}
\newcommand{\RQwenEightWrongReadbackAbstainN}{120}
\newcommand{\RQwenEightWrongReadbackAbstainRate}{0.000}
\newcommand{\RQwenEightWrongReadbackAbstainSuccess}{0.000}
\newcommand{\RQwenEightWrongReadbackAbstainWrongCommit}{0.042}
\newcommand{\RQwenEightWrongReadbackN}{120}
\newcommand{\RQwenEightWrongReadbackSuccess}{0.008}
\newcommand{\RQwenEightWrongReadbackWrongCommit}{0.692}
\newcommand{\RQwenFourCleanAlwaysAbstainRate}{0.000}
\newcommand{\RQwenFourCleanAlwaysN}{120}
\newcommand{\RQwenFourCleanAlwaysSuccess}{0.225}
\newcommand{\RQwenFourCleanAlwaysWrongCommit}{0.042}
\newcommand{\RQwenFourCleanBothAbstainRate}{0.000}
\newcommand{\RQwenFourCleanBothN}{120}
\newcommand{\RQwenFourCleanBothSuccess}{0.792}
\newcommand{\RQwenFourCleanBothWrongCommit}{0.042}
\newcommand{\RQwenFourCleanContractAbstainAbstainRate}{0.000}
\newcommand{\RQwenFourCleanContractAbstainN}{120}
\newcommand{\RQwenFourCleanContractAbstainRate}{0.000}
\newcommand{\RQwenFourCleanContractAbstainSuccess}{0.792}
\newcommand{\RQwenFourCleanContractAbstainWrongCommit}{0.042}
\newcommand{\RQwenFourCleanContractN}{120}
\newcommand{\RQwenFourCleanContractSuccess}{0.792}
\newcommand{\RQwenFourCleanContractWrongCommit}{0.042}
\newcommand{\RQwenFourCleanNeverAbstainRate}{0.000}
\newcommand{\RQwenFourCleanNeverN}{120}
\newcommand{\RQwenFourCleanNeverSuccess}{0.792}
\newcommand{\RQwenFourCleanNeverWrongCommit}{0.042}
\newcommand{\RQwenFourCleanOnErrorAbstainRate}{0.000}
\newcommand{\RQwenFourCleanOnErrorN}{120}
\newcommand{\RQwenFourCleanOnErrorSuccess}{0.758}
\newcommand{\RQwenFourCleanOnErrorWrongCommit}{0.050}
\newcommand{\RQwenFourCleanOnUncAbstainRate}{0.000}
\newcommand{\RQwenFourCleanOnUncN}{120}
\newcommand{\RQwenFourCleanOnUncSuccess}{0.750}
\newcommand{\RQwenFourCleanOnUncWrongCommit}{0.067}
\newcommand{\RQwenFourCleanReadbackAbstainAbstainRate}{0.000}
\newcommand{\RQwenFourCleanReadbackAbstainN}{120}
\newcommand{\RQwenFourCleanReadbackAbstainRate}{0.000}
\newcommand{\RQwenFourCleanReadbackAbstainSuccess}{0.792}
\newcommand{\RQwenFourCleanReadbackAbstainWrongCommit}{0.042}
\newcommand{\RQwenFourCleanReadbackN}{120}
\newcommand{\RQwenFourCleanReadbackSuccess}{0.792}
\newcommand{\RQwenFourCleanReadbackWrongCommit}{0.042}
\newcommand{\RQwenFourDelayedAlwaysAbstainRate}{0.000}
\newcommand{\RQwenFourDelayedAlwaysN}{120}
\newcommand{\RQwenFourDelayedAlwaysSuccess}{0.225}
\newcommand{\RQwenFourDelayedAlwaysWrongCommit}{0.033}
\newcommand{\RQwenFourDelayedBothAbstainRate}{0.000}
\newcommand{\RQwenFourDelayedBothN}{120}
\newcommand{\RQwenFourDelayedBothSuccess}{0.308}
\newcommand{\RQwenFourDelayedBothWrongCommit}{0.058}
\newcommand{\RQwenFourDelayedContractAbstainAbstainRate}{0.000}
\newcommand{\RQwenFourDelayedContractAbstainN}{120}
\newcommand{\RQwenFourDelayedContractAbstainRate}{0.000}
\newcommand{\RQwenFourDelayedContractAbstainSuccess}{0.792}
\newcommand{\RQwenFourDelayedContractAbstainWrongCommit}{0.042}
\newcommand{\RQwenFourDelayedContractN}{120}
\newcommand{\RQwenFourDelayedContractSuccess}{0.792}
\newcommand{\RQwenFourDelayedContractWrongCommit}{0.042}
\newcommand{\RQwenFourDelayedNeverAbstainRate}{0.000}
\newcommand{\RQwenFourDelayedNeverN}{120}
\newcommand{\RQwenFourDelayedNeverSuccess}{0.792}
\newcommand{\RQwenFourDelayedNeverWrongCommit}{0.042}
\newcommand{\RQwenFourDelayedOnErrorAbstainRate}{0.000}
\newcommand{\RQwenFourDelayedOnErrorN}{120}
\newcommand{\RQwenFourDelayedOnErrorSuccess}{0.758}
\newcommand{\RQwenFourDelayedOnErrorWrongCommit}{0.050}
\newcommand{\RQwenFourDelayedOnUncAbstainRate}{0.000}
\newcommand{\RQwenFourDelayedOnUncN}{120}
\newcommand{\RQwenFourDelayedOnUncSuccess}{0.750}
\newcommand{\RQwenFourDelayedOnUncWrongCommit}{0.067}
\newcommand{\RQwenFourDelayedReadbackAbstainAbstainRate}{0.950}
\newcommand{\RQwenFourDelayedReadbackAbstainN}{120}
\newcommand{\RQwenFourDelayedReadbackAbstainRate}{0.000}
\newcommand{\RQwenFourDelayedReadbackAbstainSuccess}{0.000}
\newcommand{\RQwenFourDelayedReadbackAbstainWrongCommit}{0.000}
\newcommand{\RQwenFourDelayedReadbackN}{120}
\newcommand{\RQwenFourDelayedReadbackSuccess}{0.308}
\newcommand{\RQwenFourDelayedReadbackWrongCommit}{0.058}
\newcommand{\RQwenFourDropAlwaysAbstainRate}{0.000}
\newcommand{\RQwenFourDropAlwaysN}{120}
\newcommand{\RQwenFourDropAlwaysSuccess}{0.000}
\newcommand{\RQwenFourDropAlwaysWrongCommit}{0.258}
\newcommand{\RQwenFourDropBothAbstainRate}{0.000}
\newcommand{\RQwenFourDropBothN}{120}
\newcommand{\RQwenFourDropBothSuccess}{0.008}
\newcommand{\RQwenFourDropBothWrongCommit}{0.392}
\newcommand{\RQwenFourDropContractAbstainAbstainRate}{0.000}
\newcommand{\RQwenFourDropContractAbstainN}{120}
\newcommand{\RQwenFourDropContractAbstainRate}{0.000}
\newcommand{\RQwenFourDropContractAbstainSuccess}{0.000}
\newcommand{\RQwenFourDropContractAbstainWrongCommit}{0.825}
\newcommand{\RQwenFourDropContractN}{120}
\newcommand{\RQwenFourDropContractSuccess}{0.000}
\newcommand{\RQwenFourDropContractWrongCommit}{0.825}
\newcommand{\RQwenFourDropNeverAbstainRate}{0.000}
\newcommand{\RQwenFourDropNeverN}{120}
\newcommand{\RQwenFourDropNeverSuccess}{0.000}
\newcommand{\RQwenFourDropNeverWrongCommit}{0.825}
\newcommand{\RQwenFourDropOnErrorAbstainRate}{0.000}
\newcommand{\RQwenFourDropOnErrorN}{120}
\newcommand{\RQwenFourDropOnErrorSuccess}{0.000}
\newcommand{\RQwenFourDropOnErrorWrongCommit}{0.808}
\newcommand{\RQwenFourDropOnUncAbstainRate}{0.000}
\newcommand{\RQwenFourDropOnUncN}{120}
\newcommand{\RQwenFourDropOnUncSuccess}{0.008}
\newcommand{\RQwenFourDropOnUncWrongCommit}{0.808}
\newcommand{\RQwenFourDropReadbackAbstainAbstainRate}{0.950}
\newcommand{\RQwenFourDropReadbackAbstainN}{120}
\newcommand{\RQwenFourDropReadbackAbstainRate}{0.000}
\newcommand{\RQwenFourDropReadbackAbstainSuccess}{0.000}
\newcommand{\RQwenFourDropReadbackAbstainWrongCommit}{0.000}
\newcommand{\RQwenFourDropReadbackN}{120}
\newcommand{\RQwenFourDropReadbackSuccess}{0.008}
\newcommand{\RQwenFourDropReadbackWrongCommit}{0.392}
\newcommand{\RQwenFourLoudAlwaysAbstainRate}{0.000}
\newcommand{\RQwenFourLoudAlwaysN}{120}
\newcommand{\RQwenFourLoudAlwaysSuccess}{0.225}
\newcommand{\RQwenFourLoudAlwaysWrongCommit}{0.033}
\newcommand{\RQwenFourLoudBothAbstainRate}{0.000}
\newcommand{\RQwenFourLoudBothN}{120}
\newcommand{\RQwenFourLoudBothSuccess}{0.425}
\newcommand{\RQwenFourLoudBothWrongCommit}{0.042}
\newcommand{\RQwenFourLoudContractAbstainAbstainRate}{0.000}
\newcommand{\RQwenFourLoudContractAbstainN}{120}
\newcommand{\RQwenFourLoudContractAbstainRate}{0.000}
\newcommand{\RQwenFourLoudContractAbstainSuccess}{0.425}
\newcommand{\RQwenFourLoudContractAbstainWrongCommit}{0.042}
\newcommand{\RQwenFourLoudContractN}{120}
\newcommand{\RQwenFourLoudContractSuccess}{0.425}
\newcommand{\RQwenFourLoudContractWrongCommit}{0.042}
\newcommand{\RQwenFourLoudNeverAbstainRate}{0.000}
\newcommand{\RQwenFourLoudNeverN}{120}
\newcommand{\RQwenFourLoudNeverSuccess}{0.425}
\newcommand{\RQwenFourLoudNeverWrongCommit}{0.042}
\newcommand{\RQwenFourLoudOnErrorAbstainRate}{0.000}
\newcommand{\RQwenFourLoudOnErrorN}{120}
\newcommand{\RQwenFourLoudOnErrorSuccess}{0.517}
\newcommand{\RQwenFourLoudOnErrorWrongCommit}{0.042}
\newcommand{\RQwenFourLoudOnUncAbstainRate}{0.000}
\newcommand{\RQwenFourLoudOnUncN}{120}
\newcommand{\RQwenFourLoudOnUncSuccess}{0.250}
\newcommand{\RQwenFourLoudOnUncWrongCommit}{0.033}
\newcommand{\RQwenFourLoudReadbackAbstainAbstainRate}{0.000}
\newcommand{\RQwenFourLoudReadbackAbstainN}{120}
\newcommand{\RQwenFourLoudReadbackAbstainRate}{0.000}
\newcommand{\RQwenFourLoudReadbackAbstainSuccess}{0.425}
\newcommand{\RQwenFourLoudReadbackAbstainWrongCommit}{0.042}
\newcommand{\RQwenFourLoudReadbackN}{120}
\newcommand{\RQwenFourLoudReadbackSuccess}{0.425}
\newcommand{\RQwenFourLoudReadbackWrongCommit}{0.042}
\newcommand{\RQwenFourMaxSteps}{15}
\newcommand{\RQwenFourModel}{qwen3:4b-q4\_K\_M}
\newcommand{\RQwenFourPartialAlwaysAbstainRate}{0.000}
\newcommand{\RQwenFourPartialAlwaysN}{40}
\newcommand{\RQwenFourPartialAlwaysSuccess}{0.000}
\newcommand{\RQwenFourPartialAlwaysWrongCommit}{0.050}
\newcommand{\RQwenFourPartialBothAbstainRate}{0.000}
\newcommand{\RQwenFourPartialBothN}{40}
\newcommand{\RQwenFourPartialBothSuccess}{0.000}
\newcommand{\RQwenFourPartialBothWrongCommit}{0.725}
\newcommand{\RQwenFourPartialContractAbstainAbstainRate}{0.000}
\newcommand{\RQwenFourPartialContractAbstainN}{40}
\newcommand{\RQwenFourPartialContractAbstainRate}{0.000}
\newcommand{\RQwenFourPartialContractAbstainSuccess}{0.000}
\newcommand{\RQwenFourPartialContractAbstainWrongCommit}{1.000}
\newcommand{\RQwenFourPartialContractN}{40}
\newcommand{\RQwenFourPartialContractSuccess}{0.000}
\newcommand{\RQwenFourPartialContractWrongCommit}{1.000}
\newcommand{\RQwenFourPartialNeverAbstainRate}{0.000}
\newcommand{\RQwenFourPartialNeverN}{40}
\newcommand{\RQwenFourPartialNeverSuccess}{0.000}
\newcommand{\RQwenFourPartialNeverWrongCommit}{1.000}
\newcommand{\RQwenFourPartialOnErrorAbstainRate}{0.000}
\newcommand{\RQwenFourPartialOnErrorN}{40}
\newcommand{\RQwenFourPartialOnErrorSuccess}{0.000}
\newcommand{\RQwenFourPartialOnErrorWrongCommit}{1.000}
\newcommand{\RQwenFourPartialOnUncAbstainRate}{0.000}
\newcommand{\RQwenFourPartialOnUncN}{40}
\newcommand{\RQwenFourPartialOnUncSuccess}{0.000}
\newcommand{\RQwenFourPartialOnUncWrongCommit}{1.000}
\newcommand{\RQwenFourPartialReadbackAbstainAbstainRate}{1.000}
\newcommand{\RQwenFourPartialReadbackAbstainN}{40}
\newcommand{\RQwenFourPartialReadbackAbstainRate}{0.000}
\newcommand{\RQwenFourPartialReadbackAbstainSuccess}{0.000}
\newcommand{\RQwenFourPartialReadbackAbstainWrongCommit}{0.000}
\newcommand{\RQwenFourPartialReadbackN}{40}
\newcommand{\RQwenFourPartialReadbackSuccess}{0.000}
\newcommand{\RQwenFourPartialReadbackWrongCommit}{0.725}
\newcommand{\RQwenFourReadAlwaysAbstainRate}{0.000}
\newcommand{\RQwenFourReadAlwaysN}{120}
\newcommand{\RQwenFourReadAlwaysSuccess}{0.133}
\newcommand{\RQwenFourReadAlwaysWrongCommit}{0.075}
\newcommand{\RQwenFourReadBothAbstainRate}{0.000}
\newcommand{\RQwenFourReadBothN}{120}
\newcommand{\RQwenFourReadBothSuccess}{0.508}
\newcommand{\RQwenFourReadBothWrongCommit}{0.158}
\newcommand{\RQwenFourReadContractAbstainAbstainRate}{0.250}
\newcommand{\RQwenFourReadContractAbstainN}{120}
\newcommand{\RQwenFourReadContractAbstainRate}{0.000}
\newcommand{\RQwenFourReadContractAbstainSuccess}{0.458}
\newcommand{\RQwenFourReadContractAbstainWrongCommit}{0.117}
\newcommand{\RQwenFourReadContractN}{120}
\newcommand{\RQwenFourReadContractSuccess}{0.508}
\newcommand{\RQwenFourReadContractWrongCommit}{0.158}
\newcommand{\RQwenFourReadNeverAbstainRate}{0.000}
\newcommand{\RQwenFourReadNeverN}{120}
\newcommand{\RQwenFourReadNeverSuccess}{0.525}
\newcommand{\RQwenFourReadNeverWrongCommit}{0.250}
\newcommand{\RQwenFourReadOnErrorAbstainRate}{0.000}
\newcommand{\RQwenFourReadOnErrorN}{120}
\newcommand{\RQwenFourReadOnErrorSuccess}{0.483}
\newcommand{\RQwenFourReadOnErrorWrongCommit}{0.225}
\newcommand{\RQwenFourReadOnUncAbstainRate}{0.000}
\newcommand{\RQwenFourReadOnUncN}{120}
\newcommand{\RQwenFourReadOnUncSuccess}{0.483}
\newcommand{\RQwenFourReadOnUncWrongCommit}{0.250}
\newcommand{\RQwenFourReadReadbackAbstainAbstainRate}{0.000}
\newcommand{\RQwenFourReadReadbackAbstainN}{120}
\newcommand{\RQwenFourReadReadbackAbstainRate}{0.000}
\newcommand{\RQwenFourReadReadbackAbstainSuccess}{0.525}
\newcommand{\RQwenFourReadReadbackAbstainWrongCommit}{0.250}
\newcommand{\RQwenFourReadReadbackN}{120}
\newcommand{\RQwenFourReadReadbackSuccess}{0.525}
\newcommand{\RQwenFourReadReadbackWrongCommit}{0.250}
\newcommand{\RQwenFourSeedFirst}{100}
\newcommand{\RQwenFourSeedSecond}{101}
\newcommand{\RQwenFourThreshold}{0.900}
\newcommand{\RQwenFourWrongAlwaysAbstainRate}{0.000}
\newcommand{\RQwenFourWrongAlwaysN}{120}
\newcommand{\RQwenFourWrongAlwaysSuccess}{0.008}
\newcommand{\RQwenFourWrongAlwaysWrongCommit}{0.250}
\newcommand{\RQwenFourWrongBothAbstainRate}{0.000}
\newcommand{\RQwenFourWrongBothN}{120}
\newcommand{\RQwenFourWrongBothSuccess}{0.058}
\newcommand{\RQwenFourWrongBothWrongCommit}{0.317}
\newcommand{\RQwenFourWrongContractAbstainAbstainRate}{0.000}
\newcommand{\RQwenFourWrongContractAbstainN}{120}
\newcommand{\RQwenFourWrongContractAbstainRate}{0.000}
\newcommand{\RQwenFourWrongContractAbstainSuccess}{0.058}
\newcommand{\RQwenFourWrongContractAbstainWrongCommit}{0.775}
\newcommand{\RQwenFourWrongContractN}{120}
\newcommand{\RQwenFourWrongContractSuccess}{0.058}
\newcommand{\RQwenFourWrongContractWrongCommit}{0.775}
\newcommand{\RQwenFourWrongNeverAbstainRate}{0.000}
\newcommand{\RQwenFourWrongNeverN}{120}
\newcommand{\RQwenFourWrongNeverSuccess}{0.058}
\newcommand{\RQwenFourWrongNeverWrongCommit}{0.775}
\newcommand{\RQwenFourWrongOnErrorAbstainRate}{0.000}
\newcommand{\RQwenFourWrongOnErrorN}{120}
\newcommand{\RQwenFourWrongOnErrorSuccess}{0.058}
\newcommand{\RQwenFourWrongOnErrorWrongCommit}{0.750}
\newcommand{\RQwenFourWrongOnUncAbstainRate}{0.000}
\newcommand{\RQwenFourWrongOnUncN}{120}
\newcommand{\RQwenFourWrongOnUncSuccess}{0.058}
\newcommand{\RQwenFourWrongOnUncWrongCommit}{0.767}
\newcommand{\RQwenFourWrongReadbackAbstainAbstainRate}{0.883}
\newcommand{\RQwenFourWrongReadbackAbstainN}{120}
\newcommand{\RQwenFourWrongReadbackAbstainRate}{0.000}
\newcommand{\RQwenFourWrongReadbackAbstainSuccess}{0.058}
\newcommand{\RQwenFourWrongReadbackAbstainWrongCommit}{0.000}
\newcommand{\RQwenFourWrongReadbackN}{120}
\newcommand{\RQwenFourWrongReadbackSuccess}{0.058}
\newcommand{\RQwenFourWrongReadbackWrongCommit}{0.308}
\newcommand{\RTasksPerSplit}{60}

\IfFileExists{authors.tex}{\input{authors}}{\author{Anonymous}}% chktex 27
\title{When Should a Tool Agent Recover? Comparing Triggers for Silent Tool Faults with Recovery Held Fixed}
\date{}
\begin{document}
\maketitle
\begin{abstract}
Tool-using agents often receive a success result from a call that did something else. Recovery
methods bundle \emph{when} to suspect a step with \emph{what} to do about it. We hold the agent and
its recovery message fixed and vary only the trigger, comparing seven triggers on matched
fault-injected tasks in a mock record service, under a pre-registered analysis. On \RQwenEightModel{},
an independent read-back of the world state beats a return-value contract on silent write faults by
\RPOneDiff{} in success (Holm $p$ \RPOneHolmP{}), almost entirely through dropped writes, which the
contract cannot see, and combining both beats never recovering by \RPThreeDiff{} (Holm $p$
\RPThreeHolmP{}). The contract does not significantly raise success on wrong reads (Holm $p$
\RPTwoHolmP{}), an inconclusive result. Exploratory runs on \RQwenFourModel{} and \RHaikuModel{}
suggest that detection and repair are separate problems: read-back flags silent write faults on every
model, but only the strongest model repairs most of them, and on the smallest the combined check
lowers success. On \RHaikuModel{}, adding the contract to read-back lowers success on wrong and partial
writes; we offer a hypothesis for why.
\end{abstract}
\section{Introduction}\label{sec:intro}
A tool-using agent usually learns that something went wrong from the tool itself: a timeout, an
error code, a malformed reply. The harder case is a tool that reports success but did something
else. A write is dropped, lands on the wrong field, or is applied only in part, and the result the
agent sees is exactly what a correct call would have returned. Agents often fail to detect such quiet
failures~\cite{kraishan2026loud}, benchmarks find implicit tool failures harder than explicit
ones~\cite{zhu2026tools}, what an agent's history says about a step can hurt when the world does not
match it~\cite{pakhomov2026cleaning}, and recent systems add monitors, contracts or verification layers
to catch such failures~\cite{panthi2026outcome,babu2026selfhealing,dip2026invariantaware,kaynar2026tools}.

These systems bundle two decisions: \emph{when} to suspect a step and \emph{what} to do about it.
When they are compared, a better trigger and a better recovery procedure cannot be told apart. We
separate them. We hold the agent, its prompt and its recovery message fixed, and vary only the trigger
that decides when the agent is told that a step may be wrong. We compare seven triggers, from
``never'' to an independent read-back of the world state after every write, on matched tasks where
every arm draws the same fault and the same fault slot.

This is a narrow controlled comparison on a mock record service, not a new method. Runtime monitors of
action postconditions already trigger replanning in classical planning~\cite{cardoso2026rvreplan},
and verification of non-atomic tool calls~\cite{mansoor2026verified}, handling of writes that may
already have taken effect~\cite{li2026where}, outcome contracts~\cite{panthi2026outcome} and state
checkers for wrong reads~\cite{zhang2026rite} exist for LLM agents. Outcome Monitors already vary when
a fixed generic warning attaches, find that the recovery-tool list carries the gain, and detect no
difference from timing~\cite{panthi2026outcome}. Our study differs in its faults, silent writes graded on the final
world state; in including an independent world-state read-back alone and combined with a return-value
check; in the full trigger set under recover, with abstain for the two checking triggers; and in a
pre-registered analysis on a small local model, with a larger hosted model as an exploratory contrast. Our contributions are:
\begin{itemize}
  \item A matched, pre-registered comparison of seven triggers and two responses across seven fault
        types, graded on the final world state (Sections~\ref{sec:method} and~\ref{sec:setup}).
  \item On Qwen3-8B, read-back beats a return-value contract on silent write faults by \RPOneDiff{} in
        success (Holm $p$ \RPOneHolmP{}), and combining both beats never recovering by
        \RPThreeDiff{} (Holm $p$ \RPThreeHolmP{}) (Section~\ref{sec:results}).
  \item An inconclusive pre-registered test: the contract does not significantly raise success on
        wrong reads (Holm $p$ \RPTwoHolmP{}).
  \item Exploratory evidence that detection and repair are separate problems: read-back flags silent
        write faults on all three models, but only the strongest repairs most of them.
\end{itemize}

\section{Related work}\label{sec:related}
\paragraph{Silent tool failures.} Benchmarks inject tool perturbations and find implicit failures
harder than explicit ones~\cite{zhu2026tools}, and agents often fail to detect results that come back
well formed but wrong~\cite{kraishan2026loud}, or, under repeated poisoning, adopt them
even after checking~\cite{anon2026tools}. AgentFaultBench checks recovery against the resulting state as well as the
agent's response~\cite{gupta2026agentfaultbench}; RITE studies silent wrong values in multi-step
e-commerce agents~\cite{zhang2026rite}. A verification-first monitor that issues corrective directives
improves recovery from tool failures on a real filesystem~\cite{kaynar2026tools}, and recoveries after
partial executions can themselves harm the external state~\cite{trofimov2026agent}. We share their
concern with faults that look like success, and add a comparison of triggers with recovery held fixed.

\paragraph{Read-back and contracts.} Runtime monitors of action preconditions and postconditions drive
replanning in classical planning~\cite{cardoso2026rvreplan}. For LLM agents, verified tool calls address
non-atomic behaviour such as delayed visibility and partial updates~\cite{mansoor2026verified}, and
LIMBO asks whether exactly-once behaviour belongs in the model, the harness or the tool
contract~\cite{li2026where}. Outcome Monitors check results against outcome contracts and issue a
receipt on a violation~\cite{panthi2026outcome}; SilentGuard enforces invariants in the
harness~\cite{dip2026invariantaware}; RECEIPT measures and repairs the assumption that a call returns
with its action finished~\cite{yang2026receipt}. Our read-back and contract triggers are simplified
versions of these ideas, and our read-back is idealised (Section~\ref{sec:method}), so that the
comparison is about the trigger and not its engineering.

\paragraph{When to intervene.} Recovery can rescue a failing trajectory or disrupt a good
one~\cite{xiao2026harnesses}, and agents often recover on their own if left alone~\cite{duan2026intervene}.
Uncertainty is a common but unproven signal for when to correct an agent~\cite{li2026uncertainty}, a
calibrated failure score need not say whether intervening helps~\cite{zhang2026calibration}, and a gate
that only replaces a proposal can act only where the base policy was already
wrong~\cite{sethi2026consistency}. Middleware can cue an
agent when it overlooks execution feedback~\cite{yang2026beyond}, and how an error is delivered changes
what the agent does next~\cite{woo2026same}. The choice between repairing, replanning and abandoning a
plan has been studied in theory~\cite{tewari2026ai}, repair has been compared with replanning
empirically~\cite{yuan2026replan}, and trajectory studies trace where errors enter and which controls
reach them~\cite{wang2026dependable}. Our \texttt{always} and \texttt{on\_uncertainty} arms are the
simplest such policies.

\paragraph{Replanning and self-correction.} Agents can replan from environment
feedback~\cite{sun2023adaplanner}, decompose sub-tasks they fail to execute~\cite{prasad2023adapt},
verbally reflect on task feedback~\cite{shinn2023reflexion} or refine their outputs with
self-feedback~\cite{madaan2023selfrefine}. Each of these acts on a failure signal; a silent fault
provides none. Without an external signal, models struggle to correct their own
reasoning~\cite{huang2023large}. Our triggers supply that external signal; the recovery itself is left
to the model.

\section{Method}\label{sec:method}
\paragraph{Service and tasks.} The agent operates a small in-memory record service with collections
of orders, items and accounts. It reads with \texttt{get\_record} and \texttt{list\_records} and
writes with \texttt{update\_field}, \texttt{transfer}, \texttt{create\_record} and
\texttt{delete\_record}. Write tools echo their effect: an update returns the record after the write
and a transfer returns both balances. Tasks come from six templates (for example, cancel an order and
restock its items, or move money between accounts), each instantiated ten times in each test split and five times in the development split. A task
succeeds only if the agent declares it done, every goal holds in the final state, and nothing else
changed. No language model judges success.

\paragraph{Agent.} Every arm runs the same loop and system prompt. At each step the model replies with
one JSON object, either a tool call with a short thought and a stated confidence (a probability), or a
declaration that the task is done; a missing confidence counts as full confidence. The tool result, or
the tool error, comes back as the next user message. A reply that is not a valid action is returned to
the model as an error; a short run of such replies in a row ends the episode as a \emph{parse
failure}. The agent has at most \RQwenEightMaxSteps{} steps. Sampling settings are fixed per model and
given in the released code.

\paragraph{Faults.} At most one fault is injected per episode. For each (task, fault, seed) a slot
$k$ is drawn and shared by every arm; the fault arms at the $k$-th successful call of its kind and
fires once, on the first eligible call at or after it. Only the slot index is shared: arms change what the agent does, so the
trajectory and whether the fault fires at all can differ between arms (Section~\ref{sec:limitations}). The faults are:
\emph{clean}, which injects nothing;
\emph{loud}, which returns an error;
three silent write faults that return the echo a correct write would return, namely \emph{drop},
which does not apply the write, \emph{wrong}, which writes the value to a different field of the
record, and \emph{partial}, which applies only the debit of a transfer or only the first half of a
new record's fields;
\emph{read}, which corrupts one field of a returned record, with an implausible value (a negative
number or an invalid order status) in a little under half of the episodes and a plausible one otherwise;
and \emph{delayed}, which applies the write but hides it from reads for two further calls. Delayed
exposes false alarms, since there is nothing to repair.

\paragraph{Triggers.} After each tool call a trigger decides whether to flag it. \texttt{never} and
\texttt{always} are the bounds; \texttt{on\_error} fires on tool errors and \texttt{on\_uncertainty}
when the agent's stated confidence is below \RQwenEightThreshold{}. \texttt{on\_contract} sees only
the call and its result and checks the schema, the field types, simple value domains (no negative numbers, a known order
status) and that the echo matches the request.
At the faulted call of a silent write fault the echo is correct, so the contract cannot flag it; it can
flag later reads of a record the fault left with a wrongly typed value or missing fields.
\texttt{on\_readback} computes the effect a write should have had from the harness's own snapshot of
the state before the call and a hand-written effect model per tool, re-reads every affected record, and
fires on any difference. The write faults change the stored state itself, so any read would show them;
read-back's advantage is that it compares what it reads with the expected effect. Its reads go through
a path that the wrong-read fault does not corrupt and that does not count towards fault slots, but
that still sees delayed visibility. This is an idealised read-back: a deployed one would need its own
expected-effect model. \texttt{readback\_contract} fires if either does.

\paragraph{Responses.} A flag injects one fixed message, whatever the trigger:
\begin{quote}\small
NOTICE\@: a check flagged a possible problem with step $k$ (\emph{tool}). Its result may not match the
real state. Decide how to continue: you may re-read state, re-issue the call, or change your plan.
\end{quote}
The agent never sees what the check found. Under \emph{recover} the agent continues with this
message; under \emph{abstain} the episode ends and counts as a failure. All seven triggers are run with
recover, and the contract and read-back triggers also with abstain, giving nine arms.

\section{Experimental setup}\label{sec:setup}
\paragraph{Models.} The confirmatory model is \RQwenEightModel{} (Qwen3-8B below), run locally through
Ollama. We also run \RQwenFourModel{} (Qwen3-4B) as a size contrast and \RHaikuModel{} (Haiku) as a
stronger model. The local models use seeds \RQwenEightSeedFirst{} and \RQwenEightSeedSecond{}; Haiku
uses one seed (\RHaikuSeedFirst{}). The final Claude run cost \RHaikuCostUsd{}; the local runs
cost nothing. The models differ in provider, family (Qwen against Claude), size, quantisation, sampling, seeds and
parse-failure rate, so ``stronger model'' below names a bundle of differences, not one factor.

\paragraph{Splits.} A development split was used for every design choice. The uncertainty threshold
was chosen there from a grid, by rescoring the confidences the agent had stated in development runs, as
the value whose flag rate on clean and delayed steps came closest to read-back's; every lower value in
the grid never fired. The test split reported here has \RTasksPerSplit{} tasks and was run once per
model with frozen configurations; a second test split of the same size was never run, is not used for
any result, and is released with the code for future work. Partial faults
need a transfer or a create, which only two templates contain, so other (task, fault) pairs are
skipped in every arm. A cell (model, fault, arm) therefore holds \RQwenEightDropReadbackN{} episodes
for the Qwen models on most faults (\RQwenEightPartialReadbackN{} on partial), and
\RHaikuDropReadbackN{} (\RHaikuPartialReadbackN{}) for Haiku.

\paragraph{Pre-registration and tests.} Before the test split was run we fixed three primary
comparisons, all on Qwen3-8B with recover: \textbf{P1}, \texttt{on\_readback} against
\texttt{on\_contract} on the three silent write faults pooled; \textbf{P2}, \texttt{on\_contract}
against \texttt{never} on wrong reads; and \textbf{P3}, \texttt{readback\_contract} against
\texttt{never} averaged over all seven faults. For each task we average the paired differences in
success (first arm minus second) over its faults and seeds; Diff.\ is the mean of these per-task
differences. Each test is a two-sided sign-flip permutation test on the per-task differences, with
Holm correction across the three. We also report a bootstrap confidence interval that resamples tasks,
and the counts of (task, fault, seed) pairs rescued (only the first arm succeeded) and harmed (only the
second did). The interval is descriptive; where it and the permutation $p$ disagree, the pre-registered
permutation test decides. Everything else, including every test on the other two models, is exploratory
and reported without correction.

\section{Results}\label{sec:results}
\begin{table}[t]
\centering\small
\begin{tabular}{@{}lrrrrr@{}}
\toprule
Test & Diff. & 95\% CI & $p$ & Holm $p$ & Rescue/harm \\
\midrule
P1 & \RPOneDiff & [\RPOneCiLo, \RPOneCiHi] & \RPOneP & \RPOneHolmP & \RPOneRescue/\RPOneHarm \\
P2 & \RPTwoDiff & [\RPTwoCiLo, \RPTwoCiHi] & \RPTwoP & \RPTwoHolmP & \RPTwoRescue/\RPTwoHarm \\
P3 & \RPThreeDiff & [\RPThreeCiLo, \RPThreeCiHi] & \RPThreeP & \RPThreeHolmP & \RPThreeRescue/\RPThreeHarm \\
\bottomrule
\end{tabular}

\caption{Pre-registered primary tests on Qwen3-8B with recover. P1: \texttt{on\_readback} minus
\texttt{on\_contract} on the three silent write faults (\RPOneNPairs{} pairs); P2:
\texttt{on\_contract} minus \texttt{never} on wrong reads (\RPTwoNPairs{} pairs); P3:
\texttt{readback\_contract} minus \texttt{never} over all seven faults (\RPThreeNPairs{} pairs). Diff.\
is the mean of per-task differences in success; the confidence interval resamples tasks; $p$ is the
sign-flip permutation test and Holm $p$ corrects it across P1--P3. Rescue/harm count (task, fault,
seed) pairs where only the first or only the second arm succeeded.}\label{tab:primary}
\end{table}

\paragraph{Primary tests.} Table~\ref{tab:primary} gives the confirmatory results. Read-back beats
the contract on silent write faults (P1: \RPOneDiff{}, interval [\RPOneCiLo{}, \RPOneCiHi{}], Holm $p$
\RPOneHolmP{}). Since the contract cannot flag these faults at the faulted call, on them it behaves
like \texttt{never}, and P1 measures what one flag on the faulted write is worth. Descriptively (the per-fault split was
pre-registered as a breakdown, not a test), almost the whole
effect comes from dropped writes (\RPOneDropDiff{}); on wrong and partial writes the differences are
\RPOneWrongDiff{} and \RPOnePartialDiff{}. P2 is not significant: the contract raises success on wrong
reads by \RPTwoDiff{} (Holm $p$ \RPTwoHolmP{}). Its interval [\RPTwoCiLo{}, \RPTwoCiHi{}] excludes
zero, so the test is inconclusive rather than evidence of no effect; the wrong-read fault fired in only
a share \RExpQwenEightAppliedContractRead{} of contract episodes, which dilutes it. Combining both checks
beats never recovering when averaged over all faults (P3: \RPThreeDiff{}, Holm $p$ \RPThreeHolmP{}), at
a cost on delayed writes, where it harms \RPThreeDelayedHarmTotal{} pairs that \texttt{never} solved.

\begin{figure}[t]
\centering
\pgfplotsset{xmin=0, xmax=1, enlarge y limits=0.06, legend columns=4}
\pgfplotsset{onlymodel/.style={y filter/.code={\edef\tempa{\thisrow{model}}%
  \ifx\tempa\tempb\else\def\pgfmathresult{inf}\fi}},
  every axis/.append style={xbar, bar width=3pt, width=0.37\textwidth, height=6.4cm,
  symbolic y coords={Never,Always,OnError,OnUnc,Contract,Readback,Both}, ytick=data,
  xlabel={success rate}, xlabel style={font=\scriptsize},
  y tick label style={font=\scriptsize}, x tick label style={font=\scriptsize}, title style={font=\small}},
  every axis plot/.append style={draw=black, line width=0.2pt}}
\newcommand{\faultbars}{%
  \addplot[fill=darkgray] table[y=arm, x=drop, onlymodel]{generated/fig_success.dat};
  \addplot[fill=gray] table[y=arm, x=wrong, onlymodel]{generated/fig_success.dat};
  \addplot[fill=lightgray] table[y=arm, x=partial, onlymodel]{generated/fig_success.dat};
  \addplot[fill=white] table[y=arm, x=read, onlymodel]{generated/fig_success.dat};}
\def\tempb{QwenFour}%
\begin{tikzpicture}\begin{axis}[title={Qwen3-4B}, legend to name=faultlegend,
  legend style={font=\small}]
\faultbars
\legend{silent drop, silent wrong, silent partial, silent read}
\end{axis}\end{tikzpicture}%
\def\tempb{QwenEight}%
\begin{tikzpicture}\begin{axis}[title={Qwen3-8B}, yticklabels={}]
\faultbars
\end{axis}\end{tikzpicture}%
\def\tempb{Haiku}%
\begin{tikzpicture}\begin{axis}[title={Haiku}, yticklabels={}]
\faultbars
\end{axis}\end{tikzpicture}

\pgfplotslegendfromname{faultlegend}
\caption{Success rate on the four silent faults for each trigger with recover. Each bar is one
(model, fault, arm) cell of \RQwenEightDropReadbackN{} episodes for the Qwen models
(\RQwenEightPartialReadbackN{} on partial) and \RHaikuDropReadbackN{} for Haiku
(\RHaikuPartialReadbackN{} on partial). Full tables for all faults and arms are in
Appendix~\ref{app:tables}.}\label{fig:success}
\end{figure}

\paragraph{Detection is not repair (exploratory).} Read-back flags silent write faults on every
model. With abstain, an episode stops at the first flag, so the abstain rate under read-back measures
how often it flags: on dropped writes \RHaikuDropReadbackAbstainAbstainRate{} for Haiku,
\RQwenEightDropReadbackAbstainAbstainRate{} for Qwen3-8B and \RQwenFourDropReadbackAbstainAbstainRate{}
for Qwen3-4B, and on wrong writes \RHaikuWrongReadbackAbstainAbstainRate{},
\RQwenEightWrongReadbackAbstainAbstainRate{} and \RQwenFourWrongReadbackAbstainAbstainRate{}. What the
agent does with the same flag under recover differs (Figure~\ref{fig:success}). On dropped writes Haiku
reaches \RHaikuDropReadbackSuccess{}, Qwen3-8B \RQwenEightDropReadbackSuccess{} and Qwen3-4B only
\RQwenFourDropReadbackSuccess{}, which is no better than \texttt{on\_uncertainty}
(\RQwenFourDropOnUncSuccess{}). On wrong writes Haiku reaches \RHaikuWrongReadbackSuccess{} and on
partial writes \RHaikuPartialReadbackSuccess{}, while Qwen3-8B reaches
\RQwenEightWrongReadbackSuccess{} and Qwen3-4B \RQwenFourWrongReadbackSuccess{} on wrong writes, the
same as its \texttt{never} arm. In the episodes we read, the small models re-apply the field they meant
to write on wrong writes and miss the corrupted sibling field; on partial transfers a retry would debit
twice. Read-back tells the agent \emph{that} a step went wrong, not \emph{what} went wrong. On
Qwen3-4B, \RExpQwenFourStatusParseFailure{} of its \RExpQwenFourEpisodes{} episodes end in a parse
failure, which depresses every arm. Its exploratory P3 difference is \RExpQwenFourPThreeDiff{}
(interval [\RExpQwenFourPThreeCiLo{}, \RExpQwenFourPThreeCiHi{}]), and \RExpQwenFourPThreeDelayedHarm{}
of its \RExpQwenFourPThreeHarm{} harmed pairs are delayed writes, where read-back's false alarms lower
success from \RQwenFourDelayedNeverSuccess{} to \RQwenFourDelayedBothSuccess{}. Uncorrected versions of P1--P3 for all models are in
Table~\ref{tab:exploratory}.

\paragraph{Recover or abstain (exploratory).} Abstaining always scores as a failure here, so recover
wins on success. What abstain buys is fewer wrong commits, episodes the agent declares done with the
state wrong: on Qwen3-8B dropped writes the wrong-commit rate is \RQwenEightDropNeverWrongCommit{} under
\texttt{never}, \RQwenEightDropReadbackWrongCommit{} under read-back with recover and
\RQwenEightDropReadbackAbstainWrongCommit{} under read-back with abstain, and on wrong writes
\RQwenEightWrongReadbackWrongCommit{} with recover against \RQwenEightWrongReadbackAbstainWrongCommit{}
with abstain. Which response is better
depends on what a wrong commit costs; the pre-registration (H5) treated this comparison as descriptive
only.

\paragraph{Wrong reads (exploratory).} The contract flags \RExpQwenEightByKindContractImplausibleFlagged{} of the
\RExpQwenEightByKindContractImplausibleN{} implausible wrong reads on Qwen3-8B and \RExpQwenEightByKindContractPlausibleFlagged{} of
the \RExpQwenEightByKindContractPlausibleN{} plausible ones. Among the implausible reads, success
rises from \RExpQwenEightByKindNeverImplausibleSuccess{} episodes under \texttt{never} to
\RExpQwenEightByKindContractImplausibleSuccess{} under the contract; on Haiku it rises from
\RExpHaikuByKindNeverImplausibleSuccess{} to \RExpHaikuByKindContractImplausibleSuccess{} of
\RExpHaikuByKindContractImplausibleN{}. On Qwen3-4B the contract flags
\RExpQwenFourByKindContractImplausibleFlagged{} of \RExpQwenFourByKindContractImplausibleN{} implausible
reads, yet success falls from \RExpQwenFourByKindNeverImplausibleSuccess{} to
\RExpQwenFourByKindContractImplausibleSuccess{}. Read-back cannot see a wrong read, and
neither checking trigger flags a plausible one; only \texttt{always} and, now and then,
\texttt{on\_uncertainty} do.

\paragraph{False alarms (exploratory).} On clean and delayed episodes read-back flags a share
\RExpQwenEightFlagReadback{} of steps on Qwen3-8B, all of them on delayed writes, while the contract
flags \RExpQwenEightFlagContract{}. \texttt{on\_error} flags \RExpQwenEightFlagOnError{}, on tool errors caused by the agent's own
calls. With recover this costs Qwen3-8B least (delayed success
\RQwenEightDelayedReadbackSuccess{} against \RQwenEightDelayedNeverSuccess{} for \texttt{never}), Haiku
more (\RHaikuDelayedReadbackSuccess{} against \RHaikuDelayedNeverSuccess{}) and Qwen3-4B most
(\RQwenFourDelayedReadbackSuccess{} against \RQwenFourDelayedNeverSuccess{}); with abstain it is fatal
(\RQwenEightDelayedReadbackAbstainSuccess{}). Flagging every step cuts both ways: on clean tasks
\texttt{always} lowers success from \RQwenFourCleanNeverSuccess{} to \RQwenFourCleanAlwaysSuccess{} on
Qwen3-4B and from \RQwenEightCleanNeverSuccess{} to \RQwenEightCleanAlwaysSuccess{} on Qwen3-8B, but on
Haiku it raises success on every silent fault, for example from \RHaikuDropNeverSuccess{} to
\RHaikuDropAlwaysSuccess{} on dropped writes. The uncertainty trigger flags a share
\RExpQwenEightFlagOnUnc{} of clean and delayed steps on Qwen3-8B, so it does fire, but its flags do
not follow the faults: success on the silent faults is about that of \texttt{never}
(Table~\ref{tab:full-qwen8}), and on loud faults it falls from \RQwenEightLoudNeverSuccess{} to
\RQwenEightLoudOnUncSuccess{}. Of the pre-registered secondary hypotheses, H3 holds descriptively on
Qwen3-8B (\texttt{readback\_contract} has the highest success over all faults of any arm), and H4 holds
only in part: \texttt{always} raises the most false flags and both it and \texttt{on\_uncertainty} end
below \texttt{never} over all faults, but \texttt{on\_uncertainty} raises fewer false flags than
\texttt{on\_error} or read-back.

\section{Discussion}\label{sec:discussion}
\paragraph{Detection and repair are separate problems (exploratory).} Apart from flagging every step, read-back is the
only trigger that flags silent write faults reliably, and it does so on all three models; the
uncertainty trigger catches a few by chance. Whether a flag turns
into success depends on the model. A dropped write is repaired by repeating the call, which Haiku
usually does and Qwen3-8B sometimes does, while Qwen3-4B almost never gets there. A wrong-value or
partial write needs the agent to find what changed and undo it, which only Haiku does often. This
matches reports that, under repeated poisoning, agents adopt wrong answers even after
checking~\cite{anon2026tools}, and that a monitor which issues corrective directives improves
recovery~\cite{kaynar2026tools}. Telling the agent \emph{what} differed from the expected state may
help the small models, as telling a client the exact fault kind does elsewhere~\cite{woo2026same}; we did not test this, to keep the NOTICE identical
across triggers.

\paragraph{More checks can hurt (exploratory).} On Haiku, adding the contract to read-back lowers
success on wrong and partial writes (\RHaikuWrongBothSuccess{} and \RHaikuPartialBothSuccess{}, against
\RHaikuWrongReadbackSuccess{} and \RHaikuPartialReadbackSuccess{} with read-back alone). These cells
hold \RHaikuWrongReadbackN{} and \RHaikuPartialReadbackN{} episodes, and every pair where only one of
the two arms succeeded favours read-back alone. An inspection of those pairs suggests a cause. Both arms
raise the same first flag, on the faulted write. With both checks, the agent's later re-reads of the
corrupted record fail the contract's type check (wrong writes) or schema check (partial writes) and
are flagged as well, and in the episodes we read the agent then ends the task without
repairing the record. Our hypothesis is that the NOTICE, which says the \emph{result} may not match the
state when in fact the \emph{state} is wrong, leads the agent to distrust the evidence it needs. This
rests on one seed and a manual inspection, and we never varied the NOTICE, so it is a hypothesis to
test. On wrong reads, by contrast, adding the contract helps Haiku (\RHaikuReadBothSuccess{} against
\RHaikuReadReadbackSuccess{}).

\paragraph{Practical implications.} These suggestions rest largely on the exploratory results. Of the
triggers here, read-back is the only one that reliably sees writes which echoed the right thing, at the
cost of extra reads and an expected-effect model, so it is worth adding where that effect can be
written down. Whether it pays depends on the model: it needs to act on true flags and to shrug off
false ones. On Qwen3-4B the combined check lowered success overall (\RExpQwenFourPThreeDiff{}), almost
all of it through false alarms on delayed writes, so where writes may become visible late, read-back
should wait or re-check before it flags. Keep recovery rather than abstaining when a wrong commit is
not much worse than an abandoned task; abstain where a wrong commit is costly. Flagging every step
hurt the small models on clean tasks and helped Haiku on silent faults, so its value depends on the
model and it is no substitute for a targeted trigger. Contracts are
cheap and catch implausible values, but they cannot see a write that echoed the right thing, and their
message should say what they found.

\section{Limitations}\label{sec:limitations}
The service is a mock with six task templates and ten instances of each per test split, so the tasks are
not independent and the effects are tied to these templates; we cluster by task but do not model the
template, and a small effect such as P2's is absent in some templates. Only
Qwen3-8B carries confirmatory tests; Haiku was run with a single seed. The read-back is idealised: it
knows the true state before each write and the expected effect of every tool (Section~\ref{sec:method}).
The implausible variant of the wrong-read fault was added during development so that the contract had a
class of faults it could see; it is reported separately. Five further issues come from our own review of
the experiment. First, the grader counts every flag that is not on the faulted call as a false flag,
so the contract's correct flags on later re-reads of a corrupted record are counted as false; the
NOTICE wording is shared by all triggers and is not tuned to any of them. Second, the faulted call and
the fault's random draws come from the same seeded generator, so they are correlated within a (task,
fault, seed) triple; every arm still uses the same slot. Third, Qwen3-4B often fails to produce a
parseable reply, which depresses all its arms. Fourth, a fault only fires if the agent reaches an
eligible call, and arms change how often it does: on Qwen3-8B the wrong-read fault fired in a share
\RExpQwenEightAppliedNeverRead{} of \texttt{never} episodes but \RExpQwenEightAppliedAlwaysRead{} of
\texttt{always} episodes. Fifth, the step budget may cap repair: \RExpQwenEightStatusBudgetExhausted{}
Qwen3-8B and \RExpQwenFourStatusBudgetExhausted{} Qwen3-4B episodes end on it, and we did not test
whether a larger budget lets the small models repair wrong and partial writes. The uncertainty
threshold was chosen by rescoring development runs, where it would have fired far less often than it
did on the test split; we report it as run rather than tune it on the test split.

\paragraph{Code and data.} The agent, service, fault injector, triggers, grader, frozen configurations,
the development split, both test splits, and the episode logs behind every number in this paper are
released at \url{https://github.com/rkumar-bengaluru/silent-tool-fault-recovery}, with a script that
regenerates all tables and figures from the logs.
\section{AI assistance}\label{sec:ai}
We used text-to-text generative AI (Anthropic's Claude, through Claude Code) to search and summarise
the literature, write and test the experiment code, run the analysis scripts and draft this paper.
Every number in the paper is generated by script from the stored results, and every citation was
checked against the cited paper's abstract or full text. The authors reviewed all of the content and take full
responsibility for it. No AI tool is listed as an author.
\label{endofmain}
\bibliographystyle{plainnat}
\bibliography{refs}
\appendix
\section{Full results}\label{app:tables}
Tables~\ref{tab:full-qwen8},~\ref{tab:full-qwen4} and~\ref{tab:full-haiku} give the success rate for
every fault and arm; bold marks the best arm per fault (ties are all bold). Each cell holds
\RQwenEightDropReadbackN{} episodes for the Qwen models (\RQwenEightPartialReadbackN{} on partial) and
\RHaikuDropReadbackN{} for Haiku (\RHaikuPartialReadbackN{} on partial). Table~\ref{tab:exploratory}
gives uncorrected versions of P1--P3 on every model. Only Table~\ref{tab:primary} is confirmatory; the
Qwen3-8B rows of Table~\ref{tab:exploratory} repeat its tests without correction.

\begin{table}[h]\centering\small
\resizebox{\textwidth}{!}{\begin{tabular}{@{}lrrrrrrrrr@{}}
\toprule
Fault & Never & Always & OnErr & OnUnc & Contr. & Readb. & Both & Contr.-Ab. & Readb.-Ab. \\
\midrule
Clean & \RQwenEightCleanNeverSuccess & \RQwenEightCleanAlwaysSuccess & \textbf{\RQwenEightCleanOnErrorSuccess} & \textbf{\RQwenEightCleanOnUncSuccess} & \RQwenEightCleanContractSuccess & \RQwenEightCleanReadbackSuccess & \RQwenEightCleanBothSuccess & \RQwenEightCleanContractAbstainSuccess & \RQwenEightCleanReadbackAbstainSuccess \\
Loud & \RQwenEightLoudNeverSuccess & \RQwenEightLoudAlwaysSuccess & \textbf{\RQwenEightLoudOnErrorSuccess} & \RQwenEightLoudOnUncSuccess & \RQwenEightLoudContractSuccess & \RQwenEightLoudReadbackSuccess & \RQwenEightLoudBothSuccess & \RQwenEightLoudContractAbstainSuccess & \RQwenEightLoudReadbackAbstainSuccess \\
Drop & \RQwenEightDropNeverSuccess & \RQwenEightDropAlwaysSuccess & \RQwenEightDropOnErrorSuccess & \RQwenEightDropOnUncSuccess & \RQwenEightDropContractSuccess & \textbf{\RQwenEightDropReadbackSuccess} & \textbf{\RQwenEightDropBothSuccess} & \RQwenEightDropContractAbstainSuccess & \RQwenEightDropReadbackAbstainSuccess \\
Wrong & \RQwenEightWrongNeverSuccess & \textbf{\RQwenEightWrongAlwaysSuccess} & \RQwenEightWrongOnErrorSuccess & \RQwenEightWrongOnUncSuccess & \RQwenEightWrongContractSuccess & \RQwenEightWrongReadbackSuccess & \RQwenEightWrongBothSuccess & \RQwenEightWrongContractAbstainSuccess & \RQwenEightWrongReadbackAbstainSuccess \\
Partial & \textbf{\RQwenEightPartialNeverSuccess} & \textbf{\RQwenEightPartialAlwaysSuccess} & \textbf{\RQwenEightPartialOnErrorSuccess} & \textbf{\RQwenEightPartialOnUncSuccess} & \textbf{\RQwenEightPartialContractSuccess} & \textbf{\RQwenEightPartialReadbackSuccess} & \textbf{\RQwenEightPartialBothSuccess} & \textbf{\RQwenEightPartialContractAbstainSuccess} & \textbf{\RQwenEightPartialReadbackAbstainSuccess} \\
Read & \RQwenEightReadNeverSuccess & \RQwenEightReadAlwaysSuccess & \RQwenEightReadOnErrorSuccess & \RQwenEightReadOnUncSuccess & \textbf{\RQwenEightReadContractSuccess} & \RQwenEightReadReadbackSuccess & \textbf{\RQwenEightReadBothSuccess} & \RQwenEightReadContractAbstainSuccess & \RQwenEightReadReadbackAbstainSuccess \\
Delayed & \RQwenEightDelayedNeverSuccess & \RQwenEightDelayedAlwaysSuccess & \textbf{\RQwenEightDelayedOnErrorSuccess} & \textbf{\RQwenEightDelayedOnUncSuccess} & \RQwenEightDelayedContractSuccess & \RQwenEightDelayedReadbackSuccess & \RQwenEightDelayedBothSuccess & \RQwenEightDelayedContractAbstainSuccess & \RQwenEightDelayedReadbackAbstainSuccess \\
\bottomrule
\end{tabular}
}
\caption{Success rate by fault and arm, \RQwenEightModel{}.}\label{tab:full-qwen8}
\end{table}
\begin{table}[h]\centering\small
\resizebox{\textwidth}{!}{\begin{tabular}{@{}lrrrrrrrrr@{}}
\toprule
Fault & Never & Always & OnErr & OnUnc & Contr. & Readb. & Both & Contr.-Ab. & Readb.-Ab. \\
\midrule
Clean & \textbf{\RQwenFourCleanNeverSuccess} & \RQwenFourCleanAlwaysSuccess & \RQwenFourCleanOnErrorSuccess & \RQwenFourCleanOnUncSuccess & \textbf{\RQwenFourCleanContractSuccess} & \textbf{\RQwenFourCleanReadbackSuccess} & \textbf{\RQwenFourCleanBothSuccess} & \textbf{\RQwenFourCleanContractAbstainSuccess} & \textbf{\RQwenFourCleanReadbackAbstainSuccess} \\
Loud & \RQwenFourLoudNeverSuccess & \RQwenFourLoudAlwaysSuccess & \textbf{\RQwenFourLoudOnErrorSuccess} & \RQwenFourLoudOnUncSuccess & \RQwenFourLoudContractSuccess & \RQwenFourLoudReadbackSuccess & \RQwenFourLoudBothSuccess & \RQwenFourLoudContractAbstainSuccess & \RQwenFourLoudReadbackAbstainSuccess \\
Drop & \RQwenFourDropNeverSuccess & \RQwenFourDropAlwaysSuccess & \RQwenFourDropOnErrorSuccess & \textbf{\RQwenFourDropOnUncSuccess} & \RQwenFourDropContractSuccess & \textbf{\RQwenFourDropReadbackSuccess} & \textbf{\RQwenFourDropBothSuccess} & \RQwenFourDropContractAbstainSuccess & \RQwenFourDropReadbackAbstainSuccess \\
Wrong & \textbf{\RQwenFourWrongNeverSuccess} & \RQwenFourWrongAlwaysSuccess & \textbf{\RQwenFourWrongOnErrorSuccess} & \textbf{\RQwenFourWrongOnUncSuccess} & \textbf{\RQwenFourWrongContractSuccess} & \textbf{\RQwenFourWrongReadbackSuccess} & \textbf{\RQwenFourWrongBothSuccess} & \textbf{\RQwenFourWrongContractAbstainSuccess} & \textbf{\RQwenFourWrongReadbackAbstainSuccess} \\
Partial & \textbf{\RQwenFourPartialNeverSuccess} & \textbf{\RQwenFourPartialAlwaysSuccess} & \textbf{\RQwenFourPartialOnErrorSuccess} & \textbf{\RQwenFourPartialOnUncSuccess} & \textbf{\RQwenFourPartialContractSuccess} & \textbf{\RQwenFourPartialReadbackSuccess} & \textbf{\RQwenFourPartialBothSuccess} & \textbf{\RQwenFourPartialContractAbstainSuccess} & \textbf{\RQwenFourPartialReadbackAbstainSuccess} \\
Read & \textbf{\RQwenFourReadNeverSuccess} & \RQwenFourReadAlwaysSuccess & \RQwenFourReadOnErrorSuccess & \RQwenFourReadOnUncSuccess & \RQwenFourReadContractSuccess & \textbf{\RQwenFourReadReadbackSuccess} & \RQwenFourReadBothSuccess & \RQwenFourReadContractAbstainSuccess & \textbf{\RQwenFourReadReadbackAbstainSuccess} \\
Delayed & \textbf{\RQwenFourDelayedNeverSuccess} & \RQwenFourDelayedAlwaysSuccess & \RQwenFourDelayedOnErrorSuccess & \RQwenFourDelayedOnUncSuccess & \textbf{\RQwenFourDelayedContractSuccess} & \RQwenFourDelayedReadbackSuccess & \RQwenFourDelayedBothSuccess & \textbf{\RQwenFourDelayedContractAbstainSuccess} & \RQwenFourDelayedReadbackAbstainSuccess \\
\bottomrule
\end{tabular}
}
\caption{Success rate by fault and arm, \RQwenFourModel{}.}\label{tab:full-qwen4}
\end{table}
\begin{table}[h]\centering\small
\resizebox{\textwidth}{!}{\begin{tabular}{@{}lrrrrrrrrr@{}}
\toprule
Fault & Never & Always & OnErr & OnUnc & Contr. & Readb. & Both & Contr.-Ab. & Readb.-Ab. \\
\midrule
Clean & \textbf{\RHaikuCleanNeverSuccess} & \RHaikuCleanAlwaysSuccess & \textbf{\RHaikuCleanOnErrorSuccess} & \textbf{\RHaikuCleanOnUncSuccess} & \textbf{\RHaikuCleanContractSuccess} & \textbf{\RHaikuCleanReadbackSuccess} & \textbf{\RHaikuCleanBothSuccess} & \textbf{\RHaikuCleanContractAbstainSuccess} & \textbf{\RHaikuCleanReadbackAbstainSuccess} \\
Loud & \textbf{\RHaikuLoudNeverSuccess} & \RHaikuLoudAlwaysSuccess & \RHaikuLoudOnErrorSuccess & \RHaikuLoudOnUncSuccess & \textbf{\RHaikuLoudContractSuccess} & \textbf{\RHaikuLoudReadbackSuccess} & \textbf{\RHaikuLoudBothSuccess} & \textbf{\RHaikuLoudContractAbstainSuccess} & \textbf{\RHaikuLoudReadbackAbstainSuccess} \\
Drop & \RHaikuDropNeverSuccess & \RHaikuDropAlwaysSuccess & \RHaikuDropOnErrorSuccess & \RHaikuDropOnUncSuccess & \RHaikuDropContractSuccess & \textbf{\RHaikuDropReadbackSuccess} & \textbf{\RHaikuDropBothSuccess} & \RHaikuDropContractAbstainSuccess & \RHaikuDropReadbackAbstainSuccess \\
Wrong & \RHaikuWrongNeverSuccess & \RHaikuWrongAlwaysSuccess & \RHaikuWrongOnErrorSuccess & \RHaikuWrongOnUncSuccess & \RHaikuWrongContractSuccess & \textbf{\RHaikuWrongReadbackSuccess} & \RHaikuWrongBothSuccess & \RHaikuWrongContractAbstainSuccess & \RHaikuWrongReadbackAbstainSuccess \\
Partial & \RHaikuPartialNeverSuccess & \RHaikuPartialAlwaysSuccess & \RHaikuPartialOnErrorSuccess & \RHaikuPartialOnUncSuccess & \RHaikuPartialContractSuccess & \textbf{\RHaikuPartialReadbackSuccess} & \RHaikuPartialBothSuccess & \RHaikuPartialContractAbstainSuccess & \RHaikuPartialReadbackAbstainSuccess \\
Read & \RHaikuReadNeverSuccess & \RHaikuReadAlwaysSuccess & \RHaikuReadOnErrorSuccess & \RHaikuReadOnUncSuccess & \textbf{\RHaikuReadContractSuccess} & \RHaikuReadReadbackSuccess & \textbf{\RHaikuReadBothSuccess} & \RHaikuReadContractAbstainSuccess & \RHaikuReadReadbackAbstainSuccess \\
Delayed & \textbf{\RHaikuDelayedNeverSuccess} & \RHaikuDelayedAlwaysSuccess & \textbf{\RHaikuDelayedOnErrorSuccess} & \textbf{\RHaikuDelayedOnUncSuccess} & \textbf{\RHaikuDelayedContractSuccess} & \RHaikuDelayedReadbackSuccess & \RHaikuDelayedBothSuccess & \textbf{\RHaikuDelayedContractAbstainSuccess} & \RHaikuDelayedReadbackAbstainSuccess \\
\bottomrule
\end{tabular}
}
\caption{Success rate by fault and arm, \RHaikuModel{}.}\label{tab:full-haiku}
\end{table}
\begin{table}[h]\centering\small
\begin{tabular}{@{}llrrrr@{}}
\toprule
Model & Test & Diff. & 95\% CI & $p$ & Rescue/harm \\
\midrule
\RHaikuModel & P1 & \RExpHaikuPOneDiff & [\RExpHaikuPOneCiLo, \RExpHaikuPOneCiHi] & \RExpHaikuPOneP & \RExpHaikuPOneRescue/\RExpHaikuPOneHarm \\
\RHaikuModel & P2 & \RExpHaikuPTwoDiff & [\RExpHaikuPTwoCiLo, \RExpHaikuPTwoCiHi] & \RExpHaikuPTwoP & \RExpHaikuPTwoRescue/\RExpHaikuPTwoHarm \\
\RHaikuModel & P3 & \RExpHaikuPThreeDiff & [\RExpHaikuPThreeCiLo, \RExpHaikuPThreeCiHi] & \RExpHaikuPThreeP & \RExpHaikuPThreeRescue/\RExpHaikuPThreeHarm \\
\RQwenFourModel & P1 & \RExpQwenFourPOneDiff & [\RExpQwenFourPOneCiLo, \RExpQwenFourPOneCiHi] & \RExpQwenFourPOneP & \RExpQwenFourPOneRescue/\RExpQwenFourPOneHarm \\
\RQwenFourModel & P2 & \RExpQwenFourPTwoDiff & [\RExpQwenFourPTwoCiLo, \RExpQwenFourPTwoCiHi] & \RExpQwenFourPTwoP & \RExpQwenFourPTwoRescue/\RExpQwenFourPTwoHarm \\
\RQwenFourModel & P3 & \RExpQwenFourPThreeDiff & [\RExpQwenFourPThreeCiLo, \RExpQwenFourPThreeCiHi] & \RExpQwenFourPThreeP & \RExpQwenFourPThreeRescue/\RExpQwenFourPThreeHarm \\
\RQwenEightModel & P1 & \RExpQwenEightPOneDiff & [\RExpQwenEightPOneCiLo, \RExpQwenEightPOneCiHi] & \RExpQwenEightPOneP & \RExpQwenEightPOneRescue/\RExpQwenEightPOneHarm \\
\RQwenEightModel & P2 & \RExpQwenEightPTwoDiff & [\RExpQwenEightPTwoCiLo, \RExpQwenEightPTwoCiHi] & \RExpQwenEightPTwoP & \RExpQwenEightPTwoRescue/\RExpQwenEightPTwoHarm \\
\RQwenEightModel & P3 & \RExpQwenEightPThreeDiff & [\RExpQwenEightPThreeCiLo, \RExpQwenEightPThreeCiHi] & \RExpQwenEightPThreeP & \RExpQwenEightPThreeRescue/\RExpQwenEightPThreeHarm \\
\bottomrule
\end{tabular}

\caption{Exploratory, uncorrected versions of P1--P3 on all three models. These $p$-values are not
corrected for multiple comparisons.}\label{tab:exploratory}
\end{table}
\end{document}
